\documentclass[11pt]{amsart}
\usepackage[margin=1.04in]{geometry}
\usepackage[T1]{fontenc}
\usepackage{lmodern,amsmath,amssymb,amsfonts,mathtools,mathrsfs,booktabs,longtable,array,microtype}
\usepackage[hidelinks]{hyperref}
\ifdefined\pdfinfoomitdate\pdfinfoomitdate=1\fi
\ifdefined\pdftrailerid\pdftrailerid{}\fi
\ifdefined\pdfsuppressptexinfo\pdfsuppressptexinfo=15\fi
\newtheorem{theorem}{Theorem}[section]
\newtheorem{proposition}[theorem]{Proposition}
\newtheorem{lemma}[theorem]{Lemma}
\newtheorem{corollary}[theorem]{Corollary}
\theoremstyle{definition}\newtheorem{definition}[theorem]{Definition}
\newtheorem{example}[theorem]{Example}
\theoremstyle{remark}\newtheorem{remark}[theorem]{Remark}
\newcommand{\E}{\mathbb E}
\newcommand{\Pp}{\mathbb P}
\newcommand{\R}{\mathbb R}
\newcommand{\N}{\mathbb N}
\newcommand{\one}{\mathbf 1}
\newcommand{\osc}{\operatorname{osc}}
\newcommand{\TV}{\operatorname{TV}}
\newcommand{\Lip}{\operatorname{Lip}}
\newcommand{\conv}{\operatorname{conv}}
\newcommand{\diag}{\operatorname{diag}}
\newcommand{\rank}{\operatorname{rank}}
\newcommand{\tr}{\operatorname{tr}}
\newcommand{\supp}{\operatorname{supp}}
\newcommand{\dist}{\operatorname{dist}}
\newcommand{\calR}{\mathcal R}
\newcommand{\calC}{\mathcal C}
\newcommand{\calE}{\mathfrak E}
\newcommand{\calG}{\mathcal G}
\newcommand{\calK}{\mathcal K}
\newcommand{\norm}[1]{\left\lVert#1\right\rVert}
\newcommand{\abs}[1]{\left|#1\right|}
\DeclareMathOperator*{\esssup}{ess\,sup}
\title[Marked physical-time transfer]{Companion F: Marked Physical-Time Transfer\\and Acquired Occupation}
\author{Qian Qi}
\date{10 October 2026}
\subjclass[2020]{Primary 60J20, 93E20; Secondary 62B15, 60K15, 94A17, 81P45}
\keywords{causal experiment, predictive quotient, acquired occupation, finite memory, Markov renewal, task-sensitive regularity, singular observation}
\begin{document}
\begin{abstract}
A finite observer may retain information across physical preparations while its boundary chain changes recurrent rank. We prove that a residual--oscillation cost for the marked Poisson equation is equivalent, up to an actual return-tail error, to the largest physical-acquisition prefix discrepancy. Its exact finite-dimensional dual gives boundary-balance lower certificates without a spectral gap. The converse retains the joint duration--exit law. For executable predictive tests it controls the actual acquired occupation, and a localized small-ball argument transfers finite-memory lower bounds at the target resolution rather than its square. Explicitly counted recursive covers give matching causal upper bounds. A predictable first-service argument derives actual geometric mass for all adaptive controllers with uniformly bounded conditional returns. A raw diagnostic family attains the general modulus against all histories for arbitrary continuous revelation profiles, including logarithmically slow acquisition. Fresh singular marks, calibration ambiguity and an attained causal defect yield a matched same-task multiresource law. Unbounded queue excursions and noncommuting quantum instruments provide independent raw verifications. Compactness, effective certification and computational resources retain their separate hypotheses.
\end{abstract}
\maketitle
\noindent\emph{Complete retained mathematical article.} The crossing convention and related notation are corrected in this copy. All substantive statements and proofs are retained. The distinct nonreset energy theorem is not used as a premise here.

\section{Introduction}
A prepared experiment determines which future tests can be executed, which histories have the same predictive content, and which distributions on that content are actually acquired. The resource problem is to turn those predictions into one causal program with finite persistent state. Formal parameter dimension is not enough: states may be reachable but scarcely acquired, a slow dynamical mode may be invisible to the task, and a good terminal quantizer need not admit an online implementation. This paper connects those distinctions through a two-sided physical-acquisition theorem.

Physical preparation need not erase the observer's evidence about an unknown parameter. With retained memory, successive physical excursions form a finite-boundary Markov-renewal process, not a sequence of independent cycles. The physical state and the reports may nevertheless be infinite. Several closed boundary classes and periodic behavior are allowed. Gain is obtained by forming a reward-to-duration ratio in each closed class and only then mixing by the actual absorption probabilities. Across a model family, the recurrent structure may change without damaging a specified task.

An approximate Poisson equation provides an upper estimate, but an upper estimate alone does not identify the resource obstruction. Let $P$ be the exit-label matrix, $T_{ij}=\E_i[\tau\one_{\{I_1=j\}}]$ the joint duration--exit matrix, $r$ the complete-excursion reward vector, and $g$ the physical-time gain. Put
\begin{equation}\label{eq:intro-poisson}
 b=r-Tg,\qquad
 k_N=\inf_u\left\{\norm{b-(I-P)u}_\infty+\frac{\osc(u)}N\right\}.
\end{equation}
The central result, Theorem~\ref{thm:modulus}, proves both that $k_N$ controls actual physical-acquisition discrepancy and that physical-acquisition discrepancies control $k_N$. The reverse comparison uses the largest prefix deviation, normalized by the total acquisition budget, and includes precisely an actual return-tail remainder. Periodic cancellation shows why the prefix cannot simply be replaced by the terminal time. A deterministic long excursion shows why boundary calculations alone cannot remove the physical-time remainder.

The converse mechanism is a finite renewal identity. Averaging the actual prefix deviations constructs a corrector whose residual is the sum of a remaining physical-tail reward and a duration-weighted terminal strip. It never separates exit labels from their durations. An exact linear-programming dual then gives signed approximate-invariance certificates for a lower bound on physical discrepancy. These certificates are not fictitious acquired distributions. Nor are they parameters or exact posterior coordinates provided to the observer. The result requires finite mean excursions uniformly over the declared entering rows, but no spectral gap, irreducibility or uniform group-inverse bound.

For executable predictive occupation tests, the same theorem becomes a two-sided temporal certificate for the actual acquired law. Theorem~\ref{thm:resolution} then localizes a small-ball witness to the intended spatial resolution $r$. A truncated squared-distance test transfers the mass whenever the occupation discrepancy is of order $r$, rather than requiring an error of order $r^2$. It retains the true mass of each lower-dimensional or singular witness. The resulting finite-acquisition lower bound combines with global candidate-valid covers and actual-law update stability to give the counted causal upper in Corollary~\ref{cor:resolution-memory}. Spatial recursive implementability remains a separately verified condition; the temporal converse does not conceal it.

Theorem~\ref{thm:profile} shows that the general modulus is attained against every history-dependent strategy on a continuum of primitive diagnostic experiments. Revelation intensity and task amplitude can be arbitrary continuous profiles. Two retained labels attain the exact minimax value at every horizon. Power profiles recover algebraic rates, while a logarithmic amplitude gives a nonalgebraic rate despite compactness and continuous gain. A queue with unbounded workload and a noncommuting quantum instrument give structurally different raw verifications. The latter crosses a preparation-rank change without a uniform recurrence bound. Neither realization is a premise of the general theorem.

For one completely specified scored experiment with fresh geometric marks, the same mechanisms give
\begin{equation}\label{eq:intro-law}
 \calR(4n,M,\delta,e)
 \asymp q_M(\rho)+\Psi_n(a,\lambda)+\delta^2
              +A_*(\varepsilon_*(e)),
\end{equation}
under explicit profile and quantization doubling. Here $\rho$ is the actually loaded mark law, $\Psi_n$ is the penalized marked modulus attained by the diagnostic task, $\delta$ is a fixed invisible calibration ambiguity, and $\varepsilon_*(e)$ is the causal deficiency of a real report-erasure interface. The exact upper and lower before doubling are also given. All four calls, controller labels, output labels and report deadlines are charged. Program description, temporary workspace and numerical precision are not identified with persistent labels.

The qualitative task-compatibility theorem is retained in Section~\ref{sec:transfer}: on a compact continuous-response family with uniformly integrable returns, continuity of gain is equivalent to uniform physical-acquisition and normalized-discount limits and uniformly finite approximate-corrector prices. Compact robust attainment and finite-model determination are consequences on a whole legal class satisfying that criterion. Finite semialgebraic instruments admit a separate effective criterion and recurrence-exhaustion schedule. Those statements do not imply unrestricted attainment when task compatibility fails.

\subsection*{Relation to earlier theory}
Multichain evaluation, classwise ratios and semi-Markov bias equations are classical~\cite{DenardoFox,FedergrunSchweitzer,LamondPuterman,Meyer}. They are not presented as new definitions of causal geometry. The literature on ergodic approximation explicitly relates convergence rates to residual--solution approximation functionals~\cite{Shaw}, and Poisson solutions generate martingale representations under established integrability conditions~\cite{GlynnInfanger}. Our contribution is the marked physical-call converse, its finite boundary-balance dual and its finite-resolution use under a fully specified causal resource model, rather than an assertion that approximation functionals or Poisson martingales are new.

Minimal predictive states and comparison of experiments likewise have substantial antecedents~\cite{ShaliziCrutchfield,Blackwell,Fritz}. The Borel realization hypotheses of our quotient lemma are explicit. Quantization and finite-memory learning are established subjects~\cite{GrafLuschgy,HellmanCover}; here actual occupation and the charged implementation are identified before these tools are used. Semialgebraic asymptotics and choice have important predecessors~\cite{BewleyKohlberg,BewleyKohlbergStationary,BasuPollackRoy}. Fixed-response linear programming and finite-algebraic certification are distinguished from efficient global synthesis. The separate primitive verification and counterexample arguments make these boundaries mathematical, rather than terminological.

\section{Experiments, programs and predictive content}
\subsection{Raw experiments and information pattern}
All physical, action, report and history spaces are standard Borel. A model parameter $\theta\in\Theta$ is fixed for the entire execution and is not an argument of the observer's program. A prepared causal experiment consists of a preparation law and normalized nonnegative marked kernels
\[
 K_\theta(dx',dy,de,d\ell\mid x,a).
\]
Here $e$ includes visible protocol flags, $\ell\in[0,1]$ is a performed score, and every transition consumes one physical call. More general positive integer costs can be expanded into marked calls when the expansion is part of the given interface. ``Positive'' means positivity preserving and normalized, not that every report has strictly positive probability. The scoring wire need not be feedback; its visibility is specified as part of the experiment. In our diagnostic examples it is never supplied to the controller or simulator. Strategies may use only prior visible reports, their current retained label, and fresh independent randomization.

An autonomous $m$-label program, $m\le M$, comprises an initial label law, action rows, report-dependent update rows and task readouts. One tuple is selected before $\theta$ is chosen. A fresh random draw used to sample a row is not retained unless encoded in the next label. There is no extra seed, history tape, exact posterior register or free internal clock. Visible protocol phases supplied by the physical experiment are distinguished from an internally stored phase. Readout numbers belong to a specified compact set; their bit representation is a separate resource. A horizon $N$ used to evaluate a stationary program is an external performance prefix, not a free internal exact-stopping clock. Any demanded internal stopping protocol must be separately represented and charged.

Physical preparations delimit excursions. Conditional on $\theta$ and the entering retained label $i$, the complete future excursion law is independent of the earlier physical history. Usually the physical preparation law is independent of $i$ and only the observer carries $i$. A prescribed label-dependent preparation is also permitted if explicitly included among the primitive legal preparation rows. This assumption is stronger than merely selecting an arbitrary return time in a hidden process. It implies that the complete conditional excursion law
\[
 \calK_x(i;d\tau,dR,dj,d\mathsf w),\qquad x=(\theta,\gamma),
\]
depends only on $i$ and $x$. Here $\tau\ge1$ is integer physical duration, $R\in[0,\tau]$ is the sum of the performed scores, $j$ is the exiting retained label, and $\mathsf w$ includes actions, reports, flags and audit marks. The observer's label is not overwritten at preparation. The triples $(\tau,R,j)$ can be arbitrarily dependent.

\subsection{Compact response families}
The following conditions will be used with a compact metric space $X$ of model--program pairs, and occasionally with an additional compact test parameter.

\begin{definition}[Continuous response family]\label{def:responses}
A family has continuous finite responses if visible finite-call cylinder probabilities and the expected bounded task readouts used below are continuous in $x$, for every entering label. This does not require total-variation continuity of a deterministic real-valued audit mark as its model value changes. It has a uniform return envelope if
\begin{equation}\label{eq:UI}
 \lim_{L\to\infty}\sup_{x,i}\E_{x,i}[\tau\one_{\{\tau>L\}}]=0.
\end{equation}
Write $\bar\mu=\sup_{x,i}\E_{x,i}\tau<\infty$ and
\begin{equation}\label{eq:tail-def}
 B_N(X)=\sup_x\frac1N\sum_{t=1}^N\max_i\E_{x,i}(\tau-t)_+.
\end{equation}
All entering rows, including rows unreachable for a particular program, occur in these suprema.
\end{definition}
Condition~\eqref{eq:UI} is not an on-policy assumption. Altering an unreachable row can change the declared program class; it is not a free normalization. It implies $B_N(X)\to0$ by a uniform tail split and Ces\`aro averaging. A useful raw sufficient condition is given next.

\begin{lemma}[Reused-row finite-response continuity]\label{lem:response}
Suppose the action and retained-label sets are finite. For each finite intervention word, the physical response has an $L^1$ density relative to fixed finite product reference measures on its distinct report coordinates, continuous in the model parameter in $L^1$. Let the program's finitely many simplex-valued report rows vary weak-star in the corresponding $L^\infty$ spaces, and let its finite action rows and compact readouts vary continuously. Assume the underlying $L^1$ spaces are separable. Then the legal closed program space is compact metrizable and finite responses are continuous. If~\eqref{eq:UI} holds, the excursion arrays $P,T,r$ defined below are continuous as well.
\end{lemma}
\begin{proof}
Positivity, row sums and any declared closed linear row restrictions are weak-star closed. Banach--Alaoglu and separability give compact metrizability of each row ball and hence of the finite product. A finite response is a finite sum of integrals of an $L^1$ density against products of bounded row entries on distinct coordinates. Approximate the density in $L^1$ by finite linear combinations of rectangle indicators. The resulting integrals are products of weak-star continuous scalar pairings. The approximation error is bounded uniformly by the $L^1$ error, even when exactly the same report row is reused at several times. There is no multiplication of weak-star limits on a single coordinate. Model continuity follows by the same uniform bound. Truncating an excursion at $L$ makes its exit probabilities, duration marks and rewards finite-response expressions. The discarded duration and reward are bounded by a constant times $\E[\tau\one_{\{\tau>L\}}]$, uniformly by~\eqref{eq:UI}. Passing $L\to\infty$ proves continuity of the arrays.
\end{proof}
Marginal domination alone is insufficient for this lemma. Finite report spaces satisfy the product condition automatically when their raw probabilities are continuous. Nondominated direct constructions will be identified separately and are not consequences of this compactness argument.

\subsection{Executable tests and the quotient}
A future test specifies a legal continuation strategy, a stopping rule within its stated budget and a bounded readout of permitted future reports or a designated audit. Histories with different available actions, protocol phases or remaining resources cannot be equated merely because some numerical predictions agree. Fix a determining countable collection $(f_j)$ of such tests, including the required interface coordinates, and their conditional predictions $q_j(h)$. Conditional versions and admissible histories are fixed once; no claim is made about arbitrary null histories under unspecified versions.

\begin{lemma}[Realized predictive quotient]\label{lem:quotient}
Suppose $q(h)=(q_j(h))_j$ is Borel, its image $S$ is standard Borel, and it admits a Borel section $s:S\to H$. Suppose equality of these coordinates determines every designated future test, including one-step continuations. Then equality of $q$ is exactly predictive equivalence for this test family. The quotient is realized by $S$ and is minimal among Borel statistics through which all the coordinates factor. If the report and continuation kernels are Borel, jointly in the action, and constant on these fibers, they descend to Borel kernels on $S$.
\end{lemma}
\begin{proof}
The determining assumption gives both implications in the equivalence. If $Z$ is another statistic and $q_j=\widetilde q_j\circ Z$ with Borel coordinate factors, then $q=\widetilde q\circ Z$ into the countable product; its restriction to $S$ gives the factorization on the realized range. Conversely any statistic factoring through $q$ cannot distinguish predictive-equivalent histories. For descended kernels evaluate the original kernel at $s(p)$. The section makes this evaluation Borel, including the action variable, and fiber invariance makes it independent of the section. The same argument applies to test evaluations and updates.
\end{proof}
Image regularity and a section are hypotheses, not consequences of positivity or of Lemma~\ref{lem:response}. A fixed-prior predictive quotient is not automatically a minimal statistic uniformly over an unknown parameter family. The resource theorems use a common predictive realization only when its model-independent initialization and updates have been verified.

\subsection{Resources and morphisms}
A causal simulator is parameter independent and includes a strategy lift, report transformation, internal state, preparation and clock map, compatible actions and a task readout. Its output must be due before each target decision that uses it. A simulator that stores $s$ labels together with an $m$-label observer uses at most $sm$ labels, with any protocol clock included in one of these coordinates. Program description, workspace, acquisition count and physical time do not follow from this label bound. A common-task deficiency compares the complete scored interface with an unchanged environment-owned audit wire; it is not unrestricted replacement of the ground-truth task. Precise finite-call and long-run error statements are given in Section~\ref{sec:composition}.

\section{The physical-acquisition modulus}\label{sec:modulus}
For each model--program pair $x$, let the conditional marked excursion have $m$ entering labels, integer duration $\tau\ge1$ and score $0\le R\le\tau$. In this section $X$ need not be compact. We assume only $\bar\mu=\sup_{x,i}\E_{x,i}\tau<\infty$, and use the actual-return quantity $B_N(X)$ of~\eqref{eq:tail-def}. Define the excursion arrays
\begin{equation}\label{eq:arrays}
 P_{ij}=\Pp_i(I_1=j),\quad T_{ij}=\E_i[\tau\one_{\{I_1=j\}}],\quad
 d=T\one,\quad D=\diag(d),\quad r_i=\E_iR.
\end{equation}
Matrices act on column rewards. Put
\begin{equation}\label{eq:time-change}
 S=I-D^{-1}(I-P),\quad c=D^{-1}r,\quad
 \Pi_A=\lim_{n\to\infty}\frac1n\sum_{j=0}^{n-1}A^j,
 \qquad g=\Pi_Sc.
\end{equation}
The stochasticity of $S$ follows from $d_i\ge1$: its diagonal is $1-(1-P_{ii})/d_i\ge0$. Every finite stochastic matrix has a semisimple eigenvalue at one and a Ces\`aro projection, regardless of periodicity. Let $J_N(x,i)$ be the expected sum of the first $N$ physical-call scores divided by $N$, from boundary label $i$, and let
\[
 J_\beta(x,i)=(1-\beta)\E_{x,i}\sum_{t\ge1}\beta^{t-1}\ell_t,
 \qquad 0<\beta<1.
\]
Write $b=r-Tg$ and $A=I-P$. Lemma~\ref{lem:marked} proves that $Pg=g$, $0\le g\le1$, and $b$ belongs to the range of $A$. In particular these definitions do not require irreducibility or aperiodicity. For $N\ge1$ set
\begin{align}
 k_N(x)&=\inf_{u\in\R^m}\left\{\norm{b_x-A_xu}_\infty+\frac{\osc(u)}N\right\},
 & K_N(X)&=\sup_{x\in X}k_N(x),\label{eq:modulus-primal}\\
 H_t(x)&=t(J_t(x)-g_x),\quad H_0=0,
 &\Delta_N(X)&=\sup_x\frac1N\max_{1\le t\le N}\norm{H_t(x)}_\infty.\label{eq:prefix-modulus}
\end{align}
The normalization in $\Delta_N$ is by $N$, not by the prefix length $t$. The vector $u$ is an analytical certificate, not a program register. The family supremum allows $u$ to depend on $x$; it does not allow the observer's program to depend on the unknown model.

\begin{theorem}[Two-sided physical-acquisition transfer]\label{thm:modulus}
For every family of the stated marked experiments with $\bar\mu<\infty$ and every integer $N\ge1$,
\begin{equation}\label{eq:modulus-comparison}
 \Delta_N(X)\le K_N(X)+B_N(X),
 \qquad K_N(X)\le(\bar\mu+2)\Delta_N(X)+B_N(X).
\end{equation}
For each $x$, the infimum in~\eqref{eq:modulus-primal} is attained. It has the exact dual
\begin{equation}\label{eq:modulus-dual}
 k_N(x)=\max\left\{z^Tb_x:
       \norm{z}_1\le1,\quad \frac12\norm{A_x^Tz}_1\le\frac1N\right\}.
\end{equation}
For each fixed $x$ and every feasible signed vector $z$, this supplies the physical-call lower certificate
\begin{equation}\label{eq:dual-lower}
 \Delta_N(X)\ge
 \frac{(z^Tb_x-B_N(X))_+}{\bar\mu+2}.
\end{equation}
If the uniform return envelope~\eqref{eq:UI} holds, then $K_N(X)\to0$ if and only if $J_N\to g$ uniformly on $X$. Neither assertion requires a uniform bound on the group inverses of $I-P_x$.
\end{theorem}
The lower certificate concerns the maximum over the displayed entering-label laws. It is not, without an accessibility argument, a lower bound from one prescribed initial distribution. The actual-occupation lower in Theorem~\ref{thm:resolution} uses its specified initial law directly, and the all-history decision lower of Theorem~\ref{thm:profile} is proved separately.

The dual vector is signed. It is not a probability, an occupation measure, or an alternative to the actual acquired law. The constraint measures its failure of invariance under the boundary kernel. For rational $P,b$ and $N$, both sides of~\eqref{eq:modulus-dual} are finite linear programs with rational data. This is a fixed-response evaluation result, not an efficient optimization theorem over Borel programs or unknown kernels.

\begin{lemma}[A physical renewal converse]\label{lem:renewal-converse}
For a fixed pair $x$, let $R^{[t]}$ be the score accumulated during the first $\min(t,\tau)$ calls of its first excursion, and define
\[
 a_t(i)=\E_i\left[R-R^{[t]}-(\tau-t)_+g(I_1)\right].
\]
For $u_N=N^{-1}\sum_{t=1}^NH_t$, one has the exact rowwise identity
\begin{equation}\label{eq:renewal-converse}
 b-(I-P)u_N
 =\frac1N\sum_{t=1}^Na_t
  +\left(\E_i\frac1N
       \sum_{s=\max(1,N-\tau+1)}^NH_s(I_1)\right)_{i=1}^m.
\end{equation}
Consequently $\norm{b-(I-P)u_N}_\infty\le B_N(X)+\bar\mu\Delta_N(X)$ and $\osc(u_N)/N\le2\Delta_N(X)$.
\end{lemma}
\begin{proof}
Both $R-R^{[t]}$ and $(\tau-t)_+g(I_1)$ lie between zero and $(\tau-t)_+$. Hence $|a_t(i)|\le\E_i(\tau-t)_+$. Conditional on the full first-excursion marks, the future starts from label $I_1$. The defining Markov-renewal property, not independence of duration and exit, gives
\[
 tJ_t(i)=\E_iR^{[t]}+
     \E_i\bigl[(t-\tau)_+g(I_1)+H_{(t-\tau)_+}(I_1)\bigr].
\]
Use $\E_i g(I_1)=g(i)$ and subtract $tg(i)$. The result is
\[
 H_t(i)=b_i-a_t(i)+\E_iH_{(t-\tau)_+}(I_1).
\]
For each realized integer $\tau\ge1$, shifting the finite sum over $t=1,\ldots,N$ leaves precisely the final $\min(\tau,N)$ terms of $\sum_{s=1}^NH_s$. Averaging and subtracting $Pu_N$ proves~\eqref{eq:renewal-converse}. The second term there is bounded by
$\E_i\min(\tau,N)\max_{s\le N}\norm{H_s}_\infty/N\le\bar\mu\Delta_N(X)$.
The first is bounded by $B_N(X)$. Each coordinate of $u_N$ is at most $N\Delta_N(X)$ in absolute value, proving its oscillation bound. All sums are finite and the only unbounded random variable used has finite mean.
\end{proof}

\begin{proof}[Proof of Theorem~\ref{thm:modulus}]
Apply Lemma~\ref{lem:physical} at every $t\le N$ with one fixed $u$, multiply by $t/N$, and use nonnegativity of the tail summands. This bounds the prefix maximum by
$\norm{b-Au}_\infty+\osc(u)/N+B_N(X)$.
Infimize separately for each pair and then take the supremum. The converse is Lemma~\ref{lem:renewal-converse} with its particular $u_N$.

Since $A\one=0$, fix $u_1=0$. A bounded objective sublevel bounds $\osc(u)$ and hence every coordinate of $u$. Continuity and finite-dimensional compactness prove attainment. To obtain the dual, write the primal as minimizing $a+(v-w)/N$ subject to
\[
 -a\one\le b-Au\le a\one,\qquad
 w\one\le u\le v\one,\qquad a\ge0.
\]
It is feasible and bounded below. The dual of the residual norm imposes $\norm{z}_1\le1$. For a zero-sum vector $y$,
\[
 \sup_{\osc(u)\le1}y^Tu=\tfrac12\norm{y}_1;
\]
indeed translate $u$ into $[0,1]^m$ and put it equal to one on the positive coordinates of $y$ and zero on the negative coordinates. Here $y=A^Tz$ has sum zero. The conjugate of $\osc(u)/N$ therefore imposes $\norm{A^Tz}_1/2\le1/N$. Finite linear-programming duality yields~\eqref{eq:modulus-dual}; its compact feasible set gives a maximum. This also proves~\eqref{eq:dual-lower}.

Finally, uniform integrability makes $B_N(X)\to0$. Uniform $J_t-g\to0$ implies $\Delta_N\to0$: split at a fixed $L$, bound the finitely many initial terms by $L/N$ using $0\le J_t,g\le1$, and bound the remaining ones by $\sup_{x,t>L}\norm{J_t-g}_\infty$. Conversely $\norm{J_N-g}\le\Delta_N$. Now use~\eqref{eq:modulus-comparison}.
\end{proof}

\begin{corollary}[Optimal balancing of approximate correctors]\label{cor:balanced-price}
With the price $\calC_X(a)$ of~\eqref{eq:corrector-price},
\begin{equation}\label{eq:balanced-price}
 K_N(X)\le\inf_{a>0}\left(a+\frac{\calC_X(a)}N\right)\le2K_N(X).
\end{equation}
Thus balancing residual tolerance and corrector oscillation loses at most a factor two relative to the pointwise penalized problem. Together with~\eqref{eq:modulus-comparison}, it is a converse as well as an upper bound for physical prefix discrepancy, up to the actual return remainder.
\end{corollary}
\begin{proof}
A corrector with residual at most $a$ gives the first inequality. Write $K=K_N(X)$. It is finite since $u=0$ gives $K\le\bar\mu$. Each pointwise minimizer has residual at most $K$ and oscillation at most $NK$. Hence $\calC_X(K+\eta)\le NK$ for every $\eta>0$; let $\eta\downarrow0$. If $K=0$, the attained zero objective forces $b_x=0$ for every $x$, and $u=0$ gives $\calC_X(a)=0$ for every $a>0$, so both sides are zero. This argument does not interchange a family supremum with an infimum over a common corrector.
\end{proof}

\begin{remark}[Why both qualifiers are needed]\label{rem:prefix-tail}
Let $P$ interchange two labels, let $\tau=1$, and let the scores at the labels be zero and one. Then $g=(1/2,1/2)^T$, $b=(-1/2,1/2)^T$, and $J_N=g$ for every even $N$, whereas $k_N=\Delta_N=1/(2N)$. Indeed $H_t=b$ for odd $t$ and zero for even $t$, so the maximal prefix norm is $1/2$ for both odd and even budgets. The primal vector $(-1/4,1/4)^T$ and dual vector $(-1/(2N),1/(2N))^T$ verify this value. A converse from the terminal discrepancy alone is false. Separately, a one-label deterministic excursion of length four, including its preparation as the first charged call, with scores $(1,1,0,0)$ has $b=0$ and $k_N=0$, but $J_1-g=1/2$. A remainder accounting for within-excursion physical time cannot be omitted. Neither example has a hidden clock supplied to the observer; the stated score patterns belong to the physical protocol.
\end{remark}

The residual--solution expression in~\eqref{eq:modulus-primal} is a finite-dimensional approximation functional. Connections between ergodic convergence rates and such functionals are classical~\cite{Shaw}; Poisson martingales and their integrability requirements are likewise established~\cite{GlynnInfanger}. The assertion here is the marked physical-time converse~\eqref{eq:renewal-converse}, its all-row tail accounting, and its use for finite-resolution acquired geometry. It does not claim the invention of approximation functionals, linear-programming duality or Poisson equations.

\section{Qualitative regularity and robust optimization}\label{sec:transfer}
For this section impose the compact continuous-response hypotheses of Definition~\ref{def:responses}, including the uniform return envelope. Use the arrays and physical-call objectives defined in Section~\ref{sec:modulus}.
For $a>0$ define the task-corrector price
\begin{equation}\label{eq:corrector-price}
 \calC_X(a)=\sup_{x\in X}\inf\left\{\osc(u):
  \norm{r-Tg-(I-P)u}_\infty\le a\right\}.
\end{equation}
The infimum is $+\infty$ if the constraint is empty; we will show it is never empty pointwise. Oscillation is $\max_i u_i-\min_i u_i$. The exact price at $a=0$ need not be uniformly finite.

\begin{theorem}[Task-visible causal transfer]\label{thm:transfer}
Under the hypotheses above, $J_N(x)\to g(x)$ and $J_\beta(x)\to g(x)$ pointwise. The following five statements are equivalent:
\begin{enumerate}
\item[(i)] The vector $g:X\to\R^m$ is continuous.
\item[(ii)] Whenever $x_j\to x$ and $\Pi_{S(x_j)}\to E$, one has $Eg(x)=g(x)$.
\item[(iii)] $\sup_x\norm{J_N(x)-g(x)}_\infty\to0$.
\item[(iv)] $\sup_x\norm{J_\beta(x)-g(x)}_\infty\to0$ as $\beta\uparrow1$.
\item[(v)] $\calC_X(a)<\infty$ for every $a>0$.
\end{enumerate}
For every $a>0$, without assuming (i), the following inequalities hold whenever their right sides are finite:
\begin{align}
 \sup_x\norm{J_N-g}_\infty
 &\le a+\frac{\calC_X(a)}{N}+B_N(X),\label{eq:physical-bound}\\
 \sup_x\norm{J_\beta-g}_\infty
 &\le a+(1-\beta)\calC_X(a)
 +(1-\beta)^2\sum_{N\ge1}N\beta^{N-1}B_N(X).\label{eq:discount-bound}
\end{align}
If $\sup_{x,i}\E_i\tau^{1+p}\le K$ for $0<p\le1$, the last tail terms are at most $mKN^{-p}$ and $mK(1-\beta)^p$, respectively. The factor $m$ can be omitted when one common stochastically dominating lifetime has this moment bound. All comparisons use the same program and the same retained-label budget.
\end{theorem}
Condition (ii) permits a change of recurrent rank. It asks only that each limiting projection fix the limiting \emph{task gain}; it does not require the projection itself to converge. Uniformity is asserted on the entire declared compact family $X$, not on arbitrary experiments lacking the response or return assumptions.

\begin{lemma}[Marked bias and physical-time gain]\label{lem:marked}
The vector $g$ is $P$-harmonic. If $C$ runs over the closed classes of $P$, $\pi_C$ is its stationary row on $C$, and $w_C(i)$ is the absorption probability from $i$, then
\begin{equation}\label{eq:class-gain}
 g_i=\sum_C w_C(i)\frac{\pi_Cr}{\pi_Cd}.
\end{equation}
Moreover $b=r-Tg$ satisfies $\Pi_Pb=0$, and $(I-P)u=b$ has a finite solution.
\end{lemma}
\begin{proof}
$S$ has the same off-diagonal graph, closed classes and absorption probabilities as $P$. On $C$ its stationary row is $\pi_C D/(\pi_Cd)$, which proves~\eqref{eq:class-gain} and harmonicity. On a closed class $g_j=g_C$ for every possible exit, so $(Tg)_i=d_i g_C$ there and $\pi_Cb=0$. Every row of $\Pi_P$ is an absorption mixture of such rows, hence $\Pi_Pb=0$. The range of $I-P$ is the kernel of $\Pi_P$. For example the group inverse
\[
 G_P=(I-P+\Pi_P)^{-1}-\Pi_P
\]
satisfies $(I-P)G_P=I-\Pi_P$, so $u=G_Pb$ is a solution. These are finite-state multichain identities, not a contraction or irreducibility assumption.
\end{proof}

\begin{lemma}[Approximate marked corrector at physical time]\label{lem:physical}
For a fixed experiment--program pair and any $u$ with $\norm{b-(I-P)u}_\infty\le a$,
\[
 \norm{J_N-g}_\infty
 \le a+\frac{\osc(u)}N+\frac1N\sum_{t=1}^N\max_i\E_i(\tau-t)_+.
\]
\end{lemma}
\begin{proof}
Let $J_0$ be the initial boundary label. Excursion $j\ge1$ enters at $J_{j-1}$ and exits at $J_j$, with duration $\tau_j$ and reward $R_j$. Put $T_0^{\rm bd}=0$, $T_j^{\rm bd}=\sum_{\ell=1}^j\tau_\ell$ and $C_N=\min\{j\ge1:T_j^{\rm bd}\ge N\}$. Then $C_N\le N$, and $\{j\le C_N\}=\{T_{j-1}^{\rm bd}<N\}$ is known before excursion $j$. In the shifted notation $S_k=T_{k-1}^{\rm bd}$ and $K_N=\min\{k\ge1:S_k\ge N\}$, one has $K_N=C_N+1\le N+1$; the number of indices $k<K_N$ is still at most $N$. All sums below use the unshifted convention.

With $e=b-(I-P)u$, the conditional mean before excursion $j$ of
\[
 R_j-\tau_jg(J_j)+u(J_j)-u(J_{j-1})
\]
is $e(J_{j-1})$. The predictable sum through $C_N$ has at most $N$ terms with finite absolute expectation. Since $Pg=g$, summation by parts gives
\[
 \sum_{j=1}^{C_N}\tau_jg(J_j)
 =T_{C_N}^{\rm bd}g(J_{C_N})-
 \sum_{j=1}^{C_N}T_{j-1}^{\rm bd}\bigl(g(J_j)-g(J_{j-1})\bigr).
\]
The second term has expectation zero by predictable bounded coefficients and bounded stopping. Also $\E g(J_{C_N})=g(J_0)$. Thus the expected completed-crossing reward minus $Ng(J_0)$ consists of a residual bounded by $aN$, a corrector boundary bounded by $\osc(u)$, and $\E[(T_{C_N}^{\rm bd}-N)g(J_{C_N})]$.

The unobserved score after physical call $N$ and this last gain term each lie between zero and the overshoot. Their difference has absolute expectation at most the expected overshoot, not twice it. There is at most one start at each integer physical time $s$, so conditioning at that predictable start gives
\[
 \E(T_{C_N}^{\rm bd}-N)\le
 \sum_{s=0}^{N-1}\max_i\E_i(\tau-(N-s))_+.
\]
Divide by $N$. No reward, duration or exit independence was introduced, and $T$ was never replaced by $DP$.
\end{proof}

\begin{proof}[Proof of Theorem~\ref{thm:transfer}]
Pointwise finite-acquisition convergence follows from Lemma~\ref{lem:physical} with the exact finite corrector of Lemma~\ref{lem:marked} and $a=0$. Its tail term vanishes by integrability. Taking the infimum over correctors and then the supremum proves~\eqref{eq:physical-bound}; an arbitrarily small excess over the infimum can be removed. For bounded scores Abel summation yields the exact identity
\begin{equation}\label{eq:abel}
 J_\beta=(1-\beta)^2\sum_{N\ge1}N\beta^{N-1}J_N.
\end{equation}
The weights sum to one, proving pointwise discount convergence and~\eqref{eq:discount-bound}.

For the moment statement,
\[
 \sum_{t=1}^N\max_i\E_i(\tau-t)_+
 \le\sum_i\E_i[\tau\min(\tau,N)]\le mK N^{1-p}.
\]
A common dominating lifetime replaces the sum over labels by one expectation. If $G$ is geometric on $\N$ with success probability $1-\beta$, Jensen's inequality gives
\[
 (1-\beta)^2\sum_{N\ge1}N^{1-p}\beta^{N-1}
 =(1-\beta)\E G^{1-p}\le(1-\beta)^p.
\]

Finite-response continuity makes each $J_N$ continuous; a uniformly convergent bounded discounted series makes $J_\beta$ continuous. Thus (iii) or (iv) implies (i). Equations~\eqref{eq:physical-bound}--\eqref{eq:discount-bound}, first with a fixed small $a$ and then with large $N$ or $\beta$ close to one, give (v)$\Rightarrow$(iii),(iv).

Suppose (i) holds. Then $b(x)=r(x)-T(x)g(x)$ is continuous and bounded, say by $B$. By Lemma~\ref{lem:marked},
\[
 A_n(x)b(x):=\frac1n\sum_{k=1}^nP(x)^k b(x)\longrightarrow0
\]
pointwise. This convergence is uniform. Indeed, for each $x_0$ choose a block length $q$ with $\norm{A_q(x_0)b(x_0)}<\epsilon$. The same bound, with $2\epsilon$, holds in a neighborhood by finite-dimensional continuity. Decomposing any $n$ into blocks of length $q$ and a remainder, and using stochastic contraction, gives $\norm{A_n b}\le2\epsilon+qB/n$ there. A finite cover of $X$ supplies a uniform bound on the chosen lengths. Now put
\[
 u_n(x)=\sum_{k=0}^{n-1}(1-k/n)P(x)^k b(x).
\]
Direct telescoping gives $(I-P)u_n=b-A_n b$ and $\osc(u_n)\le(n+1)B$. Choosing one sufficiently large $n$ for a given $a$ proves (v). No uniform bound on the exact group inverse has been used.

Finally, if $x_j\to x$ and $\Pi_{S(x_j)}\to E$, then $ES(x)=S(x)E=E$. Hence $E\Pi_{S(x)}=E$. Continuity of $c$ gives $g(x_j)\to Ec(x)=Eg(x)$ along that subsequence. Thus (i) implies (ii). Conversely the projections belong to a compact set of stochastic matrices, so every subsequence has such a further subsequence. Condition (ii) forces all resulting limits of $g(x_j)$ to equal $g(x)$, which proves (i).
\end{proof}

\begin{corollary}[Cost-universal sublevels as a special case]\label{cor:group}
If $\sup_x\norm{G_{P(x)}}_\infty\le H$, then
\[
 \sup_x\norm{J_N-g}_\infty\le\frac{2\bar\mu H}{N}+B_N(X).
\]
On a compact family of stochastic matrices, uniform boundedness of $G_P$, continuity of $\Pi_P$, locally constant recurrent rank, and uniform Ces\`aro convergence for all bounded rewards are equivalent. None is necessary for the single-task theorem.
\end{corollary}
\begin{proof}
Each coordinate of $r$ and $Tg$ lies in $[0,d_i]$, so $\norm b_\infty\le\bar\mu$. The exact solution $G_Pb$ has oscillation at most $2H\bar\mu$. For the classical equivalence use
\[
 \sum_{k=0}^{n-1}(P^k-\Pi_P)=(I-P^n)G_P.
\]
Uniform convergence implies continuity of the projection and hence locally constant integer trace. Fixed nullity at one permits a locally fixed spectral contour: otherwise an additional eigenvalue could approach one and increase nullity at the limiting stochastic matrix, whose eigenvalue one is semisimple. The contour projection and $G_P$ are then continuous; a finite cover bounds the norm. A bounded sequence of group inverses is precompact in finite dimension, and passing their defining polynomial identities to a limit shows that a norm sublevel is closed. These standard facts are recorded to delimit, not enlarge, the scope of the task theorem~\cite{Meyer,LamondPuterman}.
\end{proof}

\begin{corollary}[Robust attainment without an imposed recurrence bound]\label{cor:robust}
Let $\Theta$ and a nonempty legal program class $\Gamma_M$ be compact metric spaces, with continuous initial-label law included in the program. Assume Theorem~\ref{thm:transfer}(i) on $X=\Theta\times\Gamma_M$. Then
\begin{equation}\label{eq:robust-value}
 V_M=\min_{\gamma\in\Gamma_M}\sup_{\theta\in\Theta}g(\theta,\gamma)
 =\sup_{\substack{F\subset\Theta\\0<|F|<\infty}}
       \min_{\gamma\in\Gamma_M}\max_{\theta\in F}g(\theta,\gamma),
\end{equation}
where scalar gains use that initial law. The optimal expected $N$-call and normalized-discount values converge to $V_M$ with the errors in~\eqref{eq:physical-bound}--\eqref{eq:discount-bound}. Cluster points of asymptotically optimal programs are average optimal.
\end{corollary}
\begin{proof}
Joint gain continuity makes the worst-model objective continuous on the compact program space. For a proposed level $v$, the closed sets $\{\gamma:g(\theta,\gamma)\le v\}$ have the finite-intersection property precisely when every finite-model problem admits that level. Compactness proves~\eqref{eq:robust-value}; there is no exchange of minimization with maximization. Uniform same-program errors pass through both extrema. Compactness and continuity then prove the assertion about cluster points.
\end{proof}
When $\Gamma_M$ is the full legal class, this is an unrestricted result within the stated raw interface. When it is an architecture, the conclusion is only about that architecture. The task criterion must be proved on the whole class used in the optimization.

\section{Primitive certification and operational recurrence exhaustion}\label{sec:algebraic}
The exact criterion in Theorem~\ref{thm:transfer} does not require finite-dimensional physical dynamics. An effective consequence does. In this section all parameter and program sets are compact semialgebraic sets defined over the real algebraic numbers, all alphabets are finite, and primitive marked-response matrix entries and task readouts are continuous semialgebraic functions. Thus the finite-response continuity hypotheses remain in force. They may be rational after introducing algebraic auxiliary variables. A uniform primitive termination certificate is part of the input. An algebraic coefficient is represented by an integer polynomial and a rational isolating interval; the returned exponent is a rational number and bounds have such algebraic descriptions. Exact algebraic randomization is not automatically finite-bit physical sampling. This is not an algorithm for arbitrary computable Borel kernels.

\begin{lemma}[Finite instrumental resolvents]\label{lem:resolvent}
For a finite classical physical experiment with retained labels, let $A$ be the substochastic alive matrix on physical-state--label pairs, $B$ the one-step exit matrix to boundary labels, and $E$ the preparation matrix. Assume $\rho(A)<1$ uniformly. If $v$ is the expected one-step score from an alive pair, then
\begin{equation}\label{eq:resolvents}
 P=E(I-A)^{-1}B,\qquad
 T=E(I-A)^{-2}B,\qquad
 r=E(I-A)^{-1}v.
\end{equation}
The same formulas hold for finite-dimensional quantum instruments after replacing $A$ by the completely positive alive superoperator on the direct sum of Hermitian physical blocks indexed by observer labels, and replacing matrix rows by preparation states and exit/score functionals. In either case the excursion arrays are semialgebraic functions of the primitive data.
\end{lemma}
\begin{proof}
An excursion exiting after $j\ge1$ calls contributes $EA^{j-1}B$ to $P$ and $jEA^{j-1}B$ to $T$. The absolutely convergent sums are $(I-A)^{-1}$ and $(I-A)^{-2}$; the reward sum is the alive resolvent applied to the one-step score. In the quantum case work with unnormalized states. Composition and summation of marked completely positive maps are linear on the finite real vector space of Hermitian blocks. An exit functional applied to $A^{j-1}$ has exactly the same duration weight $j$. A uniform trace-decay certificate gives absolute convergence. Inversion of a finite nonsingular real matrix is rational in its entries. No normalization at a zero-probability report or closure of an unbounded operator is involved.
\end{proof}
A fixed preparation call can be included as an extra transient block; it must not be removed from $T$ or $r$. A positive classical or quantum Lyapunov functional giving $\norm{A^j}\le C\alpha^j$, $\alpha<1$, is a sufficient uniform termination input and supplies every finite return moment.

\begin{theorem}[Decidable task regularity and algebraic modulus]\label{thm:algebraic}
For the finite algebraic instruments just described, gain continuity on a specified compact semialgebraic family $X$ is decidable. If it holds, there exist effectively obtainable constants $a_0>0$, $C<\infty$ and rational $\kappa\ge0$ such that
\begin{equation}\label{eq:algebraic-price}
 \calC_X(a)\le C a^{-\kappa},\qquad 0<a<a_0.
\end{equation}
Consequently the acquisition error is
\begin{equation}\label{eq:algebraic-rate}
 O\bigl(N^{-1/(\kappa+1)}\bigr)+B_N(X),
\end{equation}
and the discount error is $O((1-\beta)^{1/(\kappa+1)})$ plus its tail term. The optimal robust average value and whether it is attained on the declared semialgebraic program class are decidable even when gain continuity fails. These conclusions are decision and global certification results, not polynomial-time synthesis claims.
\end{theorem}
\begin{proof}
By Lemma~\ref{lem:resolvent}, $P,T,r,S,c$ are semialgebraic. The ergodic projection is semialgebraic on each rank stratum. One direct formula is the coefficient of $z^{-1}$ in $(zI+I-S)^{-1}$: if the determinant has first nonzero term $a_qz^q$, semisimplicity at zero of $I-S$ makes that coefficient the coefficient of $z^{q-1}$ in the adjugate divided by $a_q$. There are only finitely many possible $q$. Thus the graph of $g=\Pi_Sc$ is a finite semialgebraic union. Continuity can be expressed by the first-order condition
\[
 \forall x\in X\ \forall\epsilon>0\ \exists\eta>0\
 \forall y\in X:\quad \norm{x-y}<\eta\Longrightarrow
 \norm{g(x)-g(y)}<\epsilon.
\]
Quantifier elimination over real closed fields decides it.

For the price, impose $u_1=0$, introduce upper and lower bounds for its coordinates, and encode the residual inequality in~\eqref{eq:corrector-price}. The infimum of the resulting finite linear-inequality problem and its supremum over $X$ have semialgebraic graphs, by quantifier elimination; an infimum need not be selected to make this statement. Under continuity they are finite for every $a>0$ by Theorem~\ref{thm:transfer}. A finite one-variable semialgebraic function has polynomial growth at an endpoint. The semialgebraic growth theorem, together with effective real-algebraic decomposition, supplies~\eqref{eq:algebraic-price} with rational exponent and algebraic bounds~\cite{BasuPollackRoy}. Balancing $a$ with $Ca^{-\kappa}/N$ proves~\eqref{eq:algebraic-rate}; the same balance works for discount.

Finally, the graph of the worst-model gain is semialgebraic using quantified upper bounds and arbitrarily close lower witnesses. Its infimum on the program set is therefore a real algebraic number when all defining coefficients are algebraic. The sentence asserting a program with that value decides attainment. The argument applies to a discontinuous objective as well, but then asserts neither compact attainment nor uniform finite-acquisition convergence unless the corresponding criterion holds. The algebraic assertions use standard real-algebraic machinery, not a numerical sample of the parameter continuum.
\end{proof}
The exponent in~\eqref{eq:algebraic-price} is a certified possible exponent, not necessarily the sharp exponent. The example in Section~\ref{sec:sharpness} identifies the exact price order directly. Algebraic descriptions of row probabilities are finite programs for real-arithmetic sampling; conversion to a particular bit-level random sampler, its variable running time and its approximation errors are separate implementation obligations.

\subsection{An actionable exhaustion beyond compact attainment}
A different conclusion remains available when the average objective is discontinuous. It must not be confused with Theorem~\ref{thm:algebraic}'s compatible case.

\begin{theorem}[Polynomial near-optimal recurrence schedule]\label{thm:schedule}
Assume the finite algebraic setting above, a uniform bound $\bar\mu$ on mean physical duration and a uniform return envelope over all legal programs. Suppose also that for each fixed program $\gamma$, the boundary support graph of $P(\theta,\gamma)$ is independent of $\theta\in\Theta$. Let
\[
 V=\inf_{\gamma\in\Gamma_M}\sup_{\theta\in\Theta}g(\theta,\gamma).
\]
There are $a_0>0$, $C<\infty$, rational $\kappa\ge0$ and a semialgebraic choice of common programs $\gamma_a$, $0<a<a_0$, such that
\begin{equation}\label{eq:schedule}
 \sup_\theta g(\theta,\gamma_a)\le V+a,
 \qquad
 \sup_\theta\norm{G_{P(\theta,\gamma_a)}}_\infty\le Ca^{-\kappa}.
\end{equation}
In particular the same program $\gamma_a$ satisfies, for every $N$,
\begin{equation}\label{eq:schedule-physical}
 \sup_\theta J_N(\theta,\gamma_a)
 \le V+a+2\bar\mu Ca^{-\kappa}/N+B_N(X),
\end{equation}
with the analogous discount bound. The data determine whether the unrestricted average infimum is attained. Without task compatibility, no assertion is made that the unrestricted finite-$N$ optimal values converge to $V$.
\end{theorem}
\begin{proof}
For fixed $\gamma$, the common graph fixes the number of closed classes as $\theta$ varies. The spectral-contour argument in Corollary~\ref{cor:group}, applied on compact $\Theta$, makes
\[
 H(\gamma)=\max_\theta\norm{G_{P(\theta,\gamma)}}_\infty
\]
finite. Across programs this function can be unbounded and discontinuous, but it is semialgebraic by the finite rank-stratum formulas and quantifier elimination. The set of pairs $(a,\gamma)$ satisfying the first inequality of~\eqref{eq:schedule} is nonempty for every $a>0$ and semialgebraic. Semialgebraic choice supplies a curve $\gamma_a$. The finite semialgebraic function $H(\gamma_a)$ has polynomial growth as $a\downarrow0$, proving the second inequality after shrinking $a_0$. Corollary~\ref{cor:group} gives~\eqref{eq:schedule-physical}. Attainment is the first-order statement proved decidable in Theorem~\ref{thm:algebraic}.
\end{proof}
This theorem converts recurrence exhaustion into a primitive-data, near-optimal implementation schedule. It permits the required recurrence level to diverge and does not manufacture an optimizer at the limit. The common-support assumption is testable in the finite algebraic setting; it is not a consequence of continuity alone. In particular, the rank-changing rare-herald family below uses Theorem~\ref{thm:transfer}, not this common-support theorem. Semialgebraic strategy asymptotics in stochastic games have classical precedents~\cite{BewleyKohlberg,BewleyKohlbergStationary}; the schedule here applies them to one fixed-memory robust causal program and its separately charged physical-return bridge.

\section{Actually acquired occupation and causal resolution}\label{sec:geometry}
The task criterion is useful only if the relevant geometry is itself acquired by the experiment. Fix a declared exploration and a common predictive realization when one is available. At designated scored calls let $p_t=\E[Y_t\mid H_t]$ be the conditional mean of a bounded finite-dimensional target, with $p_t\in K$ for a compact $K\subset\R^q$. In a robust family this formula is interpreted modelwise, and a common executable representation must be verified separately. Define the excursion occupation measures
\[
 \nu_i(A)=\E_i\sum_{t\text{ scored in excursion}}\one_{\{p_t\in A\}}.
\]
Preparation or other zero-score calls are omitted from this numerator but included in the duration $d_i$. The actual physical-time measure from initial label $i$ is
\begin{equation}\label{eq:occupation}
 \nu_\infty^i=\sum_C w_C(i)\frac{\sum_{j\in C}\pi_C(j)\nu_j}{\pi_Cd}.
\end{equation}
All quantization identities below concern a specified prepared law. In a robust family, modelwise quantization gives a lower bound, whereas an implementable upper must use one common codebook and recursion. The measure in~\eqref{eq:occupation} is a subprobability measure. Transient boundary occupation has zero asymptotic mass. Let $\nu_N^i$ be the expected empirical scored occupation divided by $N$ physical calls, and define $\nu_\beta^i$ with normalized-discount weights. These measures are neither preparation measures nor freely chosen volume measures.

\begin{theorem}[Task-visible acquired-law criterion]\label{thm:occupation}
Assume continuous finite responses and~\eqref{eq:UI}. Let the scored feature maps be such that finite responses of every bounded Lipschitz function of $p_t$ are continuous in $x$. Then the following are equivalent: continuity of $x\mapsto\nu_\infty^i$ in the bounded-Lipschitz metric for all initial labels; uniform convergence $\nu_N^i\to\nu_\infty^i$ in that metric; uniform discounted convergence; and Theorem~\ref{thm:transfer}'s task criterion on the compact family of normalized nonnegative Lipschitz tests of $K$. Its quantitative bounds use the supremum of the corresponding approximate-corrector prices.
\end{theorem}
\begin{proof}
Use $\mathcal F=\{f:K\to[0,1]:\Lip(f)\le1\}$, which is compact in the uniform norm by Arzel\`a--Ascoli. Treat $f$ as an additional parameter and score $f(p_t)$ only at designated calls. Formula~\eqref{eq:class-gain} is exactly~\eqref{eq:occupation} tested against $f$. Finite-response continuity is joint in $(x,f)$ because uniform changes in $f$ change normalized expected scores by at most their uniform norm. Continuity of the measure implies joint continuity of the gain; conversely continuity for all $f$ on compact $X\times\mathcal F$ gives bounded-Lipschitz continuity. Apply Theorem~\ref{thm:transfer} on this product. This argument retains the true mass of every stratum and does not require a smooth coordinate chart.
\end{proof}

For a finite positive measure $\nu$ on $K$, define
\begin{equation}\label{eq:quantization}
 q_J(\nu)=\inf_{c_1,\ldots,c_J\in\conv K}
       \int\min_{1\le j\le J}\norm{p-c_j}^2\,d\nu(p).
\end{equation}
In finite-dimensional compact support the infimum is attained. The equality-of-infima formulation remains meaningful in a general Hilbert space, where attainment requires an additional argument. Put $D_K=\operatorname{diam}(K)$.

\begin{lemma}[Projection, actual mass and quantization]\label{lem:quantization}
For a fixed exploration, the best checkpoint risk with $J$ fixed response values is its conditional-variance baseline plus $q_J$ of the actual scored occupation. Every online observer with at most $J$ such values has at least this risk, even when those labels are also used for control. Further,
\begin{equation}\label{eq:q-BL}
 |q_J(\nu)-q_J(\nu')|
 \le(D_K^2+2D_K)\,d_{\rm BL}(\nu,\nu'),
\end{equation}
with a metric using bounded-plus-Lipschitz norm at most one. If $\lambda\le\nu$ has mass $w>0$ and
\[
 \lambda(B(z,r))\le C r^d\qquad(z\in\conv K,\ 0<r\le r_0),
\]
then, whenever $r=(w/(2CJ))^{1/d}\le r_0$,
\begin{equation}\label{eq:smallball}
 q_J(\nu)\ge\frac w2\left(\frac{w}{2CJ}\right)^{2/d}.
\end{equation}
A global cover of $K$ by $J$ radius-$h_J$ balls gives $q_J(\nu)\le\nu(K)h_J^2$.
\end{lemma}
\begin{proof}
Conditional squared projection gives
$\E\norm{Y-c}^2=\E\norm{Y-p}^2+\E\norm{p-c}^2$ for each measurable response, with independent randomized choices included by conditioning. For a fixed codebook the pointwise nearest response is optimal. A checkpoint can make this assignment measurably; any online response is bounded below by the same pointwise distance. The function $\min_j\norm{p-c_j}^2$ has bound $D_K^2$ and Lipschitz constant at most $2D_K$. Taking infima proves~\eqref{eq:q-BL}. The union of $J$ radius-$r$ balls has $\lambda$-mass at most $JCr^d=w/2$; the remaining mass contributes at least $wr^2/2$. The cover bound is immediate. No mass is normalized away.
\end{proof}
The baseline may vary with the exploration. The pointwise lower remains valid when the observer jointly controls and predicts, but it is then the occupation law generated by \emph{that} observer that enters.

\subsection{A counted causal implementation}
A geometric cover alone is not an executable predictor. The following assumptions specify the missing causal content. Fix an exploration using $Q$ retained control labels. The predictive realization may use a state carrier larger than $K$: write its state as $s$, and require that the conditional task mean is $p=f(s)\in K$. Suppose it has a common known initialization at physical preparation, a jointly Borel update $F(s,a,y)$ and a task map $f$ with Lipschitz constant $L_f$. When the carrier is the conditional-mean space itself, take $f$ to be the identity. All quantization measures below are measures of $p=f(s)$, not of an unrelated carrier coordinate. A budget-independent relation between exact and candidate states is specified. For every $J$ there is a globally valid $J$-center projection of error at most $h_J$ on the full candidate domain; the update, projection and numerical perturbations preserve that relation. The report kernel is continuous on that domain in the declared task or total-variation modulus. Continuity is not a substitute for the following same-report stability estimate.

Along the \emph{actual} reports and exploration actions, suppose every admissible candidate recursion satisfies
\begin{equation}\label{eq:gain-recursion}
 e_t\le L_t e_{t-1}+h_J+\eta_{\rm num},
 \qquad
 G_t=1+L_tG_{t-1},\quad G_{\rm prep}=1,
 \qquad e_{\rm prep}\le h_J+\eta_{\rm num}.
\end{equation}
The predictor, but not the retained control label, is initialized at preparation. Assume
\begin{equation}\label{eq:occupied-stability}
 \sup_{x,N}\frac1N\E_x\sum_{t\le N,\ t\text{ scored}}G_t^2\le S_0<\infty.
\end{equation}
A limiting or horizon-dependent version gives the corresponding limiting or horizon-dependent bound. This is an actual-acquisition condition, not stability on a formal reachable set.

\begin{theorem}[Acquired geometry to causal memory]\label{thm:geometry}
Under these assumptions, a $QJ$-label observer implements the exploration and a predictor whose squared excess risk at every physical horizon is at most
\begin{equation}\label{eq:geometric-upper}
 \left[L_f(h_J+\eta_{\rm num})\sqrt{S_0}+\xi\right]^2,
\end{equation}
where $\xi$ bounds the numerical task readout error. Its retained control coordinate survives preparation. The checkpoint excess is exactly $q_J(\nu_N)$; every online $M$-label observer has excess at least $q_M(\nu_N)$ for its own actual law. Under Theorem~\ref{thm:occupation}, finite-acquisition and average quantization laws differ by~\eqref{eq:q-BL} and the task-corrector modulus.

For optimized average risk, suppose every legal $M$-label observer has baseline at least $B_*$ and its own actual occupation contains the same uniform small-ball witness parameters $(w,C,d,r_0)$. Suppose a fixed $Q$-label exploration attains $B_*$ and satisfies the global cover and stability assumptions with $h_J\le C_1J^{-1/d}$. Then, for $M\ge2Q$ and exact arithmetic,
\begin{equation}\label{eq:optimized-geometry}
 cM^{-2/d}\le\inf_{\gamma\in\Gamma_M}\calR_\infty(\gamma)-B_*
 \le C_2M^{-2/d}.
\end{equation}
The lower quantifiers are over all legal observers; the upper uses the displayed exploration.
\end{theorem}
\begin{proof}
Store $(q,c)$, the exploration label and the predictive center. Actions use only $q$, so the physical exploration law is unchanged. On an actual report, update the candidate from $c$ and project it; the next control label is updated by its original rule. There are exactly $QJ$ available pairs, including any internal clock already charged to the exploration. At preparation reset only the predictor coordinate to its common approximate prior. Induction in~\eqref{eq:gain-recursion} gives $e_t\le(h_J+\eta_{\rm num})G_t$. The task map and numerical readout yield pointwise root error at most $L_f(h_J+\eta_{\rm num})G_t+\xi$. Minkowski's inequality in the actual normalized scored occupation and~\eqref{eq:occupied-stability} give~\eqref{eq:geometric-upper}; its square, not its root, is the excess risk. Lemma~\ref{lem:quantization} and Theorem~\ref{thm:occupation} prove the checkpoint and temporal assertions.

For optimized matching, apply~\eqref{eq:smallball} to the law of each competing observer, using its at most $M$ task readouts even if the labels also implement control. This gives the uniform lower above $B_*$. For the upper choose $J=\lfloor M/Q\rfloor$, which is at least $M/(2Q)$, and use the attaining exploration with the counted product implementation. These are the extra assumptions required to pass from fixed-exploration geometry to optimized control; no such passage is automatic.
\end{proof}

\begin{proposition}[Survival-weighted expansion]\label{prop:survival}
Suppose an excursion survives each further scored step with conditional probability at most $s<1$, and every candidate-valid one-step Lipschitz multiplier is at most $L$. If $\sqrt{s}L<1$, then an actual-law stability constant as in~\eqref{eq:occupied-stability} is finite, bounded by a constant depending only on $s,L$ and the fixed number of non-scored preparation calls. Individual updates may expand, $L>1$.
\end{proposition}
\begin{proof}
At depth $k$, $G_k\le\sum_{j=0}^kL^j$. Minkowski in the weighted sequence space gives
\[
 \left(\sum_{k\ge0}s^kG_k^2\right)^{1/2}
 \le\sum_{j\ge0}L^j\left(\sum_{k\ge j}s^k\right)^{1/2}
 =\frac1{\sqrt{1-s}(1-\sqrt{s}L)}.
\]
Conditioning at each actual preparation and using at most one preparation per physical call bounds normalized cumulative occupation by this excursion sum, up to the fixed convention at preparation. No conditioning on successful survival removes its probability. Expected logarithmic contraction or random-product hypotheses may be weaker, but are not asserted here.
\end{proof}
This sufficient stability mechanism and the task-visible occupation criterion address different obstructions. The former realizes a recursive predictor; the latter permits temporal rank degeneration invisible to its test family. Neither implies the other without the displayed hypotheses.

\section{Acquisition at the target geometric resolution}\label{sec:resolution}
Let $K\subset\R^q$ be compact, with diameter at most one after fixing the units of the squared task. Put
\[
 \mathcal F_K=\{f:K\to[0,1]:\Lip(f)\le1\},\qquad
 d_K(\nu,\nu')=\sup_{f\in\mathcal F_K}\left|\int f\,d(\nu-\nu')\right|.
\]
The constant one belongs to $\mathcal F_K$, so $d_K$ controls total scored mass as well as location. The measures below are the actual subprobability occupation measures of Section~\ref{sec:geometry}, with physical preparation and unscored calls still charged in the denominator. We require that the predictive feature used on an excursion is an executable Borel function of that excursion and its entering boundary label, with a prescribed predictive initialization at preparation. More generally it suffices that its conditional excursion law has this property. An unrecorded posterior about a nonreset latent quantity may not be inserted into $K$ for free. If such information is required, the boundary state or the common predictive realization must first be supplied and accounted for.

Apply Theorem~\ref{thm:modulus} to the family $(x,f)\in X\times\mathcal F_K$, with score $f(p)$ on a scored call and zero otherwise. Denote the resulting modulus by $K_N^{K}$ and set $\zeta_N=K_N^{K}+B_N(X\times\mathcal F_K)$. Conditional excursion scores lie in $[0,\tau]$, so the theorem applies without changing the physical return envelope. In particular, for every entering label,
\begin{equation}\label{eq:occupation-prefix}
 \sup_x\max_{t\le N}\frac tN
       d_K(\nu_t^x,\nu_\infty^x)\le\zeta_N.
\end{equation}
There is also the converse of Theorem~\ref{thm:modulus}, with the same prefix supremum and constant $\bar\mu+2$. The compactness of $\mathcal F_K$ is needed only for the qualitative compact-family theorem, not for these quantitative inequalities.

\begin{theorem}[Resolution-localized acquired mass]\label{thm:resolution}
Suppose a submeasure $\lambda\le\nu_\infty$ has mass $w>0$ and
\[
 \lambda(B(z,r))\le A(r)
 \quad(z\in\conv K,
       \ 0<r\le r_0).
\]
For any $M\ge1$ and $0<r\le\min(1,r_0)$ with $M A(r)\le w/2$,
\begin{equation}\label{eq:local-mass}
 q_M(\nu_N)\ge\frac w2r^2-2r\,d_K(\nu_N,\nu_\infty)
              \ge\frac w2r^2-2r\zeta_N.
\end{equation}
Thus $\zeta_N\le wr/8$ implies $q_M(\nu_N)\ge wr^2/4$. This conclusion retains the true mass $w$; it does not normalize the submeasure to a probability. It holds for arbitrary concentration profiles $A$, including singular and stratified occupation, not just a prescribed power law.
\end{theorem}
\begin{proof}
Fix any $M$ centers $C\subset\conv K$ and put
$\varphi(p)=\min\{\dist(p,C)^2,r^2\}$. The map $\dist(p,C)$ is one-Lipschitz, and the truncated square is $2r$-Lipschitz. Moreover $0\le\varphi\le r^2$, and $r\le1$ gives $0\le\varphi/(2r)\le r/2\le1$, so $\varphi/(2r)\in\mathcal F_K$. At least $w-MA(r)\ge w/2$ of the submeasure lies at distance at least $r$ from $C$. Therefore
\[
 \int\dist(p,C)^2\,d\nu_N
 \ge\int\varphi\,d\nu_N
 \ge\frac w2r^2-2r\,d_K(\nu_N,\nu_\infty).
\]
Infimize over $C$. Equation~\eqref{eq:occupation-prefix} at $t=N$ proves the remaining assertions. Centers can be restricted to $\conv K$ by Euclidean projection; compactness gives attainment if needed. No density with respect to ambient volume has been assumed.
\end{proof}

\begin{corollary}[Finite-acquisition causal memory matching]\label{cor:resolution-memory}
Consider one fixed bounded squared prediction task and a class of legal observers with at most $M$ retained labels and at most $M$ fixed responses for this scored task (including any scored phase in that interface). Different control labels may share a readout; there are at most $M$ distinct codepoints. Assume the following statements uniformly over every competing observer, under that observer's own executed law: its actual limiting predictive occupation has a submeasure as in Theorem~\ref{thm:resolution} with common $w,A$; its full-history baseline is at least a specified $B_*(N)$; and its test-family prefix modulus is bounded by the declared $\zeta_{N,M}$. Assume also that a fixed $Q$-label exploration attains $B_*(N)$ and has the common predictive realization, global relation-preserving $J$-covers of error $h_J$, report continuity and actual occupation-weighted update stability of Theorem~\ref{thm:geometry}.

For $M\ge2Q$ choose $0<r_M\le\min(1,r_0)$ with $M A(r_M)\le w/2$ and $\zeta_{N,M}\le wr_M/8$. If $h_{\lfloor M/Q\rfloor}\le C_0r_M$, then
\begin{equation}\label{eq:resolution-memory}
 \frac w4r_M^2\le
 \inf_{\gamma:\,|\gamma|\le M}\bigl(\calR_N(\gamma)-B_*(N)\bigr)
 \le\bigl[L_f\sqrt{S_0}(C_0r_M+\eta_{\rm num})+\xi\bigr]^2.
\end{equation}
Here $\eta_{\rm num}$ is the certified per-update numerical or calibration error in the preserved relation, and $\xi$ is the task-readout error. The upper program uses $Q\lfloor M/Q\rfloor\le M$ labels. For the fixed exploration, the checkpoint excess is exactly its acquired $q_M$, and hence has the same order whenever the matching cover and mass conditions hold.
\end{corollary}
\begin{proof}
Condition on the full legal observation history for each competing program. The conditional-mean projection identity splits its risk into its own full-history baseline and the squared error of its finite readout. Randomized readouts cannot reduce that squared error below the best of their at most $M$ state-conditional means. Thus the excess over $B_*(N)$ is at least its own $q_M(\nu_N)$; Theorem~\ref{thm:resolution} applies separately to that law. The counted recursive construction of Theorem~\ref{thm:geometry}, with $J=\lfloor M/Q\rfloor$, supplies the upper bound and all program, clock and initialization conditions. It compares with the same finite-$N$ baseline by assumption; an asymptotic gain is not substituted for that baseline. The checkpoint statement is the same projection identity followed by a nearest-center encoder.
\end{proof}
For $A(r)=Cr^d$, one may take $r_M=(w/(2CM))^{1/d}$ provided $M\ge w/(2C\min(1,r_0)^d)$, as well as the other displayed budget restrictions. If $\zeta_{N,M}\le A_0N^{-\alpha}$ uniformly in the memory range, it suffices that
\[
 N\ge \left(\frac{8A_0}{wr_M}\right)^{1/\alpha},
 \qquad\text{of order }M^{1/(\alpha d)}.
\]
This is a finite-resolution acquisition condition, not an unconditional assertion that every optimal-control problem has that sample complexity. Without a uniform all-competitor witness or an attaining exploration, only the fixed-exploration and individual-observer inequalities hold.

\begin{proposition}[The resolution scale cannot be ignored]\label{prop:resolution-sharp}
Let $\nu$ be uniform measure on $[0,1]$, and let $\nu_M$ put mass $1/M$ on each midpoint $(j-1/2)/M$. Then
\[
 d_K(\nu,\nu_M)\le\frac1{4M},\qquad
 q_M(\nu_M)=0,\qquad q_M(\nu)=\frac1{12M^2}.
\]
Thus a perturbation of order the cell radius can destroy the entire finite-memory lower bound, even in one dimension.
\end{proposition}
\begin{proof}
Couple a uniform point with the midpoint of its length-$1/M$ interval. The expected distance is $1/(4M)$, which bounds every one-Lipschitz test difference. The atomic quantization assertion is immediate. For the uniform measure, a Voronoi interval of length $l$ contributes at least $l^3/12$, attained by its midpoint. Convexity gives $\sum_{j\le M}l_j^3\ge1/M^2$, with equality for equal intervals. The measures also have a direct prepared-experiment interpretation: prepare $V$ uniformly, reveal either $V$ or just its interval label, and perform the squared task against the same hidden $V$. The posterior means have laws $\nu$ and $\nu_M$, respectively. A continuous report is not a free finite-bit representation.
\end{proof}
The upper transport estimate for arbitrary squared distortion would instead lose a constant times $d_K$ independent of the cell radius. The truncation in~\eqref{eq:local-mass} is what retains the factor $r$. Both the small-ball quantization argument and its elementary one-dimensional sharpness calculation are standard in origin~\cite{GrafLuschgy}; the new use here is the marked physical-acquisition certificate at that resolution.

\section{A sharp retained-information experiment}\label{sec:sharpness}
Consider a fixed unknown $\theta\in[-1,1]$, put $t=|\theta|$ and $s=\operatorname{sign}(\theta)$ when $t>0$. One diagnostic call requires a decision $u\in\{-1,+1\}$ before its report. Its score is $t\one_{\{u\ne s\}}$, with score zero at $t=0$. The score is recorded on an environment-owned audit wire and is not supplied as feedback, now or later. The visible report is $s$ with probability $\lambda(t)$ and $\bot$ otherwise, independently across physical preparations conditional on $\theta$. Here $\lambda$ is continuous, $\lambda(0)=0$, and $\lambda(t)>0$ for $t>0$. The raw probabilities and scores are continuous on $[-1,1]$; the vanishing factor $t$ removes the apparent sign discontinuity at zero. The observer's label is retained.

\begin{theorem}[Exact diagnostic value and sharp corrector rate]\label{thm:diagnostic}
For every $n\ge1$ and every $M\ge2$, even allowing arbitrary full-history competing policies, the minimax $n$-call risk is
\begin{equation}\label{eq:diagnostic-value}
 D_n(\lambda)=\sup_{0\le t\le1}\frac{t}{2n}
                   \sum_{j=0}^{n-1}(1-\lambda(t))^j.
\end{equation}
One stationary two-label program attains it for all $n$: initialize its sign label fairly, retain it on $\bot$, and replace it by every conclusive report. Its normalized-discount value is
\begin{equation}\label{eq:diagnostic-discount}
 D_\beta(\lambda)=\sup_{0\le t\le1}
 \frac{t(1-\beta)}{2[1-\beta(1-\lambda(t))]}.
\end{equation}
For $\lambda(t)=t^k$ with integer $k\ge2$, these values are respectively of orders $n^{-1/k}$ and $(1-\beta)^{1/k}$. The two-label program satisfies Theorem~\ref{thm:transfer} on the entire parameter interval, although its group-inverse norm and its exact task bias are unbounded. Its approximate-corrector price satisfies
\begin{equation}\label{eq:sharp-price}
 \calC(a)\asymp_k a^{-(k-1)}\qquad(a\downarrow0).
\end{equation}
A stationary one-label program has minimax average risk $1/2$ when $\lambda(1)=1$; two retained labels have minimax average risk zero.
\end{theorem}
\begin{proof}
Fix $t>0$ and average the two models $\theta=t,-t$ with equal weights, solely for a minimax lower bound. Until the first conclusive report their observable histories have the same law. Their decisions may depend on arbitrary past actions and private fresh randomization, but the conditional average error probability before that report is $1/2$. At decision $j+1$, this event has probability $(1-\lambda(t))^j$. Therefore every policy has worst-model risk at least the expression in~\eqref{eq:diagnostic-value}, for each $t$. The two-label program has exactly this risk under each sign: the initial fair guess remains fair before revelation and is correct afterwards. It attains the lower for every $t$ simultaneously, proving the optimized claim. Summing the discounted geometric series gives~\eqref{eq:diagnostic-discount}.

For $\lambda(t)=t^k$,
\[
 \frac{t}{2n}\sum_{j<n}(1-t^k)^j
 \le\frac12\min(t,t^{1-k}/n).
\]
Splitting at $t=n^{-1/k}$ gives the upper rate. At that value of $t$, the sum divided by $n$ is bounded below by a positive absolute constant (with $n=1$ handled directly), giving the lower. The discount bounds follow by splitting at $t=(1-\beta)^{1/k}$; the denominator is $(1-\beta)+\beta t^k$, and $\beta\le1/2$ is absorbed by a fixed constant.

For positive $\theta$, order the labels as correct, incorrect. Then
\[
 P_t=\begin{pmatrix}1&0\\t^k&1-t^k\end{pmatrix},\quad
 r_t=\binom{0}{t},\quad \tau=1,\quad g_t=0.
\]
For negative $\theta$ exchange the labels; at zero $P_0=I$ and $r_0=0$. Thus the gain is identically zero from either entering label while the recurrent rank changes. The group inverse has norm of order $t^{-k}$ and an exact corrector needs oscillation $t^{1-k}$. An approximate corrector with residual at most $a$ must have oscillation at least $(t-a)_+/t^k$. Conversely choose an oscillation $(t-a)_+/t^k$ between the incorrect and correct labels. This attains that residual bound. Hence
\[
 \calC(a)=\sup_{0<t\le1}(t-a)_+/t^k,
\]
which is comparable to $a^{1-k}$ by the choices $t=2a$ for the lower and $t\ge a$ for the upper. Theorem~\ref{thm:transfer} now produces the sharp rate by balancing its two terms.

A one-label stationary program has the same decision distribution before every report and cannot retain whether a revelation occurred. At $\theta=\pm1$ its worst error is at least $1/2$, attained by fair decisions. The two-label program has average zero for each parameter, and the uniform bound already proved makes its worst-model finite-horizon risk tend to zero. No internal horizon or uncharged cycle number is used.
\end{proof}
This is a task-sensitive retained-learning law, not a claim about ordinary bandit feedback with observed losses. Finite-memory learning itself is classical~\cite{HellmanCover}; the role of the example here is to attain both the uniform corrector modulus and the all-policy causal acquisition value through a recurrent-rank change.

\begin{example}[Duration and exit cannot be separated]\label{ex:marked}
There are three boundary labels. From label 1, exit to label 2 with probability $1/2$ after one zero-score call, or to label 3 with probability $1/2$ after three zero-score calls. Labels 2 and 3 are absorbing with unit-duration cycles and respective scores zero and one. Then
\[
 g=(1/2,0,1)^\top,\quad d=(2,1,1)^\top,\quad
 T_{1\cdot}=(0,1/2,3/2).
\]
The first coordinate of $r-Tg$ is $-3/2$, whereas that of $r-Dg$ is $-1$. From label 1 the exact finite-call risk is $1/2-3/(2N)$ for $N\ge3$. Thus the duration--exit mark matters even in a three-state experiment. This is a check on the information pattern of the bias equation, not a new counterexample to classical semi-Markov theory.
\end{example}

\begin{example}[What the task criterion excludes]\label{ex:discontinuity}
Replace the rare-revelation loss amplitude $t$ by one on the incorrect label while keeping $P_t$ above and assigning the same positive incorrect-label loss at $t=0$. From that label the average is zero for $t>0$ and one at $t=0$. Finite responses remain continuous in this version because the reward no longer depends on $t$, but gain continuity fails. Approximate corrector prices are infinite for every $a<1$, and no uniform acquisition limit is asserted. Conversely constant zero rewards make every boundary degeneration invisible. These two cases show why cost-universal recurrence and task-visible regularity are genuinely different conditions.
\end{example}

\section{Two raw realizations}\label{sec:realizations}
\subsection{Controlled queue excursions}
A physical preparation places one job in a queue. At every service call the workload $W\ge1$ moves to $W+1$ with probability $p_\theta(a)$ and to $W-1$ otherwise, ending the excursion on hitting zero. Actions and finite noisy workload reports may depend on the retained observer state and past visible reports. The preparation call is charged and does not overwrite the observer label. Scores are bounded in $[0,1]$. Assume compact parameter sets, continuous raw report probabilities and action probabilities, and
\begin{equation}\label{eq:queue-drift}
 0\le p_\theta(a)\le\bar p<1/2
\end{equation}
for every legal action. Finite report kernels can be arbitrary functions of workload subject to this finite-response continuity; they do not alter the killed workload drift.

Fix a directed graph on the $m$ observer labels. Updates are allowed only along its reflexive reachability relation. On a first-service-call end, every designated graph edge has conditional update probability at least $\zeta>0$. The required row floors must be feasible: $k_i\zeta\le1$ for every row with $k_i$ required edges. These are a compact support-stable program architecture, not all possible queue controllers.

\begin{theorem}[Unbounded physical queue realization]\label{thm:queue}
The queue architecture above satisfies the continuous-response and uniform-return hypotheses of Theorem~\ref{thm:transfer}. It has a primitive, program-uniform group-inverse bound and therefore satisfies the task criterion for every bounded continuous scored task. Robust attainment, finite-model determination and the same-memory physical-call and discount bounds hold on this entire architecture. A certified approximation can truncate the physical workload using the explicit exponential tail below. No finite-dimensional filter or geometry of the full workload posterior is assumed.
\end{theorem}
\begin{proof}
Choose $c=\frac12\log((1-\bar p)/\bar p)$ when $\bar p>0$; the case $\bar p=0$ has deterministic immediate service termination. With $V(w)=e^{cw}$, the killed transition satisfies, uniformly over history-adaptive actions,
\[
 \E[V(W_{n+1})\one_{\{W_{n+1}>0\}}\mid W_n=w,a]
 \le\rho V(w),\qquad \rho=2\sqrt{\bar p(1-\bar p)}<1.
\]
Induction gives $\Pp(T_0>n)\le e^c\rho^n$ for the service hitting time from one. Adding the single preparation call preserves a common exponential return envelope, independent of entering label and controller. Finite-call workloads are bounded by the call count plus one, so only finitely many raw probabilities enter each response. Their continuity and Lemma~\ref{lem:response} give finite-response continuity.

The first service call ends with probability at least $1-\bar p$. Consequently each designated edge of the boundary matrix has probability at least $a_*= (1-\bar p)\zeta$. The boundary graph has exactly the same terminal strongly connected components as the designated graph: it contains every designated edge, and added edges can only follow an already existing directed path. Put $L=\max(1,m-1)$. From any state there is a path of at most $L$ edges to one selected root in a reachable terminal component. In each block of $L$ boundary transitions, the probability of hitting the set of roots is at least $a_*^L$. Thus the maximal expected boundary count to that set is at most $H_*=L/a_*^L$.

For completeness, a root-hitting bound gives $\norm{G_P}_\infty\le4H_*$. For a bounded reward $f$, center it by $b=f-\Pi_Pf$ and stop its sum at the selected root. Its expected absolute sum is at most $2H_*\norm f_\infty$. In a recurrent component the stationary centering makes the mean sum over a root-to-root return zero, so this stopped sum solves $(I-P)v=b$ also at the roots; transient states satisfy the equation by first-step conditioning. Subtracting $\Pi_Pv$ identifies $G_Pf$ and gives the factor four. Corollaries~\ref{cor:group} and~\ref{cor:robust} now apply. The bound is conservative and can be replaced by direct low-dimensional evaluation when available. Truncation errors are controlled by the common exponential return tail, not by discarding rare busy periods at zero cost.
\end{proof}
This realization verifies a genuinely unbounded physical experiment. It does not prove a full predictive-dimension law for the queue workload. A separate finite geometric task can be attached and analyzed by Theorem~\ref{thm:geometry}, but that task must not be mislabeled as the entire workload quotient. Nor does this architecture theorem claim an unrestricted solution of average-cost queue control.

\subsection{A noncommuting rank-changing quantum instrument}
Let the physical Hilbert space be $\mathbb C v\oplus\mathbb C^2$, with $v$ orthogonal to the qubit sector. Choose unit qubit vectors $\psi_+,\psi_-$ with real overlap $c_0\in(0,1)$. The unknown fixed parameter is $\theta\in[-1,1]$ and the preparation is
\begin{equation}\label{eq:quantum-prep}
 \rho_\theta=(1-t^k)|v\rangle\langle v|
               +t^k|\psi_s\rangle\langle\psi_s|,
 \qquad t=|\theta|,
\end{equation}
with $\rho_0=|v\rangle\langle v|$. Each preparation is a real physical call. The next call performs the fixed three-outcome instrument with effects
\begin{equation}\label{eq:USD}
 E_+=\frac{|\psi_-^\perp\rangle\langle\psi_-^\perp|}{1+c_0},\qquad
 E_-=\frac{|\psi_+^\perp\rangle\langle\psi_+^\perp|}{1+c_0},\qquad
 E_\bot=I-E_+-E_-.
\end{equation}
The first two effects vanish on the vacuum sector. One may use the Kraus maps $\sqrt{E_z}\rho\sqrt{E_z}$; the next preparation is fresh. The diagnostic decision and score are as in Section~\ref{sec:sharpness}, with the decision due before the instrument report. The physical phase is a visible protocol type, not an internal counter supplied to the observer.

\begin{theorem}[Noncommuting singular realization]\label{thm:quantum}
This raw instrument realizes the diagnostic theorem with $\lambda(t)=(1-c_0)t^k$. Each excursion charges one preparation and one instrument call (the diagnostic decision is due before that report), hence has duration two, retain the observer label and satisfy the task-visible criterion through $\theta=0$, where the preparation rank and recurrent rank change. It has no uniform boundary group-inverse bound. The exact all-history diagnostic value per physical call at $N=2n$ is $D_n(\lambda)/2$, and the same two-label program attains it. The physical-time rate is $N^{-1/k}$. The construction uses noncommuting physical states and a classical finite observer; it is not a theorem about arbitrary quantum memory or arbitrary additional measurements.
\end{theorem}
\begin{proof}
Write $\psi_\pm=\cos a\,e_0\pm\sin a\,e_1$, with $0<a<\pi/4$ and $c_0=\cos(2a)$. In the qubit sector $E_++E_-=\diag(\tan^2 a,1)\le I$, proving positivity and normalization. There is no false conclusive outcome, and the correct conclusive probability on $\psi_s$ is
\[
 \frac{1-c_0^2}{1+c_0}=1-c_0.
\]
The vacuum always produces $\bot$. Hence the complete visible kernel has exactly the claimed revelation probability. For $t>0$ the two state preparations corresponding to opposite signs do not commute: their qubit rank-one projections have overlap strictly between zero and one. For $0<t<1$ the preparation has rank two, at $t=0$ rank one, and at $t=1$ it is again rank one. These are actual state and support changes, not formal coordinate singularities.

All raw matrix entries and scored responses are continuous at zero. The report alphabet is finite, the physical duration is exactly two, and the two-label gain is identically zero from both entering labels. Thus Theorem~\ref{thm:transfer} applies without a recurrence sublevel. The calculation in Theorem~\ref{thm:diagnostic}, with the constant factor $1-c_0$, gives both the minimax formula and the corrector rate. Every preparation and diagnostic measurement is charged, so $N=2n$ and the score per physical call is half the per-diagnostic score. Odd horizons differ by at most one bounded call. The standard unambiguous-discrimination mechanism~\cite{ClarkeCheflesBarnettRiis} is not claimed as new; it provides an explicit noncommuting primitive realization of the task-visible theorem.
\end{proof}
With $\cos a=\sqrt3/2$ and $\sin a=1/2$, all primitive coefficients are algebraic and Theorem~\ref{thm:algebraic} also applies. In contrast, Theorem~\ref{thm:schedule}'s common-support hypothesis fails at zero. This difference is intentional: task compatibility is the mechanism that passes the singular endpoint.

\begin{corollary}[Raw verification of the acquisition modulus]\label{cor:raw-modulus}
The controlled queue architecture of Theorem~\ref{thm:queue} satisfies Theorem~\ref{thm:modulus} for every bounded task with its conditional marked excursion property. With the graph certificate stated there, both $K_N$ and $B_N(X)$ are $O(N^{-1})$, uniformly over the declared models and programs.

In Theorem~\ref{thm:quantum}, replace the excited-sector preparation weight $t^k$ by any continuous $\lambda(t)$ as in Theorem~\ref{thm:profile}, and the wrong-sign audit amplitude $t$ by $a(t)$. The same fixed noncommuting effects have conclusive probability $c_*\lambda(t)$, where $c_*=1-c_0>0$. The physical excursion has duration two, the retained observer still has two labels, and
\[
 k_N(t)=a(t)\min\{1,(Nc_*\lambda(t))^{-1}\},\qquad B_N(X)=1/N.
\]
At $N=2n$ calls its minimax risk per physical call is $D_n(a,c_*\lambda)/2$, attained against all legal histories at the same retained-state budget. This family verifies the general theorem through the rank change, with no imposed uniform bound on the exact task corrector.
\end{corollary}
\begin{proof}
The queue proof supplies an action-uniform exponential tail for physical return and a graph-floor bound for the finite retained boundary group inverse. Thus $\sup\E\tau^2<\infty$ and $B_N(X)\le\sup\E\tau^2/N$ using the common dominating lifetime. Its exact marked solution has uniformly bounded oscillation by Corollary~\ref{cor:group}, so $K_N\le C/N$. This argument starts from the service and report kernels and says nothing about an unproved finite dimension of the full workload posterior.

For the quantum family, the matrices and effect computation in Theorem~\ref{thm:quantum} depend on the excited-sector weight only through a scalar multiplier. Substituting $\lambda(t)$ preserves positivity, normalization, no false conclusive report, and finite-response continuity. The common vacuum at zero has zero task amplitude. Noncommutation of the two excited preparations at $t>0$ and their rank changes are unchanged. Preparation and measurement are distinct charged calls, giving $T=2P$, $r=(0,a(t))^T$ and $g=0$. The same primal and dual vectors used in Theorem~\ref{thm:profile} prove the displayed $k_N$. The deterministic-duration tail is exactly $1/N$. The inconclusive-history lower proof uses the actual measurement report probability $c_*\lambda(t)$ and still applies when its maximum is less than one; only the erasure-deficiency endpoint argument of Theorem~\ref{thm:profile-joint} would then require $e$ in the attainable intensity range. No new quantum memory or unlisted instrument is available to the observer.
\end{proof}

\subsection{Full geometric verification with endogenous physical returns}
The following raw subclass verifies the geometric hypotheses as well as the temporal ones. Its target is a declared task factor, not a claim that a workload or a quantum-state quotient has the same dimension as that factor.

At each physical preparation draw a fresh $V\sim\rho$ on compact $K\subset\R^q$, independently of past observer information. The preparation call is charged and the mark is hidden. A second, unscored load call reveals $V$. At least one subsequent service call is compulsory. At every service call a fixed state readout $c\in\conv K$ is scored against $V$ \emph{before} the next service report; $V$ is not supplied again. The score is an environment-owned audit, not feedback. Service ends at a visible flag and the next preparation starts a new mark. Let $L\ge1$ be the number of service calls. Physical state, reports and end hazards may depend on $V$, on the model, on actions and on the entire allowed causal history. Assume, conditionally on every entering history, every $V=v$ and every allowed adaptive strategy,
\begin{equation}\label{eq:load-mean}
 \E[L\mid\text{entering history},V=v]\le L_*,\qquad L_*<\infty.
\end{equation}
The complete physical excursion has $\tau=2+L$, so put $\mu_*=2+L_*$. The condition includes unused observer rows. The mark is a current report only on the load call; any use on later calls must pass through retained state.

\begin{theorem}[Finite-acquisition geometry under controlled returns]\label{thm:controlled-load}
For the stated raw class and every retained observer, its actual expected scored occupation through $N$ physical calls satisfies the measure inequalities
\begin{equation}\label{eq:load-occupation}
 \left(\frac1{\mu_*}-\frac1N\right)_+\rho
       \ \le\ \nu_N\ \le\ L_*\rho.
\end{equation}
For stationary finite-label observers the same bounds, with the lower coefficient $1/\mu_*$, hold for $\nu_\infty$. This limit uses the actual classwise duration weights when there are multiple boundary classes.

Suppose $\rho$ is known and common to all models, the legal class includes one fixed service action, and every observer has at most $M$ fixed service readouts. For every $N\ge2\mu_*$ and $M\ge1$ the robust optimal scored risk per physical call obeys
\begin{equation}\label{eq:load-risk}
 \frac{q_M(\rho)}{2\mu_*}
 \le\inf_{|\gamma|\le M}\sup_\theta\calR_N^\theta(\gamma)
 \le L_*q_M(\rho).
\end{equation}
A nearest-center load followed by an unchanged held center attains the upper with $M$ labels and a parameter-independent fixed service action. The conditional full-history baseline is zero. Global $M$-covers of $K$ and any actual small-ball witness for $\rho$ therefore give matching checkpoint and online geometric rates, including singular supports. No boundary recurrence constraint is imposed on competing observers.
\end{theorem}
\begin{proof}
Use the boundary convention of Lemma~\ref{lem:physical}: excursion $j\ge1$ starts at $T_{j-1}^{\rm bd}$, ends at $T_j^{\rm bd}$ and $C_u=\min\{j\ge1:T_j^{\rm bd}\ge u\}$ for $u\ge1$; set $C_0=0$. There are $C_u\le u$ starts strictly before $u$. The indicators $\one_{T_{j-1}^{\rm bd}<u}$ are predictable before the fresh mark. Conditional expectation gives
\[
 u\le\E T_{C_u}^{\rm bd}
 =\sum_{j=1}^{u}\E[\one_{T_{j-1}^{\rm bd}<u}\tau_j]
 \le\mu_*\E C_u.
\]
For $N\ge3$, exactly the starts at integer times $T_{j-1}^{\rm bd}<N-2$ guarantee that the preparation, load and first scored service are all within the first $N$ physical calls. Each such start is selected before its mark is drawn, and the first service is compulsory. Consequently
\begin{equation}\label{eq:load-first-count}
 \nu_N(A)\ge\frac{\rho(A)}N\E C_{N-2}
 \ge\frac{N-2}{\mu_*N}\rho(A).
\end{equation}
Since $\mu_*\ge3$, this implies $(1/\mu_*-1/N)_+\rho(A)$. For $N\le2$ the positive part is zero. This argument uses predictable guaranteed starts, rather than subtracting an unweighted final mark on an event that can depend on that mark.

For the upper, include complete service contributions from every started excursion. Conditioning before each fresh mark, its expected contribution to $A$ is at most $L_*\rho(A)$ by~\eqref{eq:load-mean}. There are at most $N$ starts, proving the upper inequality. For a stationary finite-label program, finite conditional means give the classwise renewal occupation formula. In each row the expected service occupation is between $\rho$ and $L_*\rho$, and its mean duration is at most $\mu_*$. Dividing separately in each recurrent class and mixing proves the asserted limiting bounds. Existence also follows from the finite-class marked reward argument for bounded tests.

For the finite-horizon risk lower it is safer to use the same guaranteed first-service calls directly, since a general observer may change its prediction during later service. Conditional on the pre-preparation history, the fresh $V$ has law $\rho$ and must be encoded into at most $M$ fixed readouts before its first service score. Its mean squared loss is at least $q_M(\rho)$, even with randomization and control information in those labels. The predictable guaranteed-start argument~\eqref{eq:load-first-count} therefore gives at least $(1/\mu_*-1/N)_+q_M(\rho)$ per physical call. All models obey this inequality. For the upper, load a nearest-center index and retain it until the next load. Conditional on $V=v$, the expected complete service loss is at most $L_*\dist(v,C_M)^2$. Including service scores after the prefix can only enlarge this nonnegative loss. Summing at most $N$ started excursions therefore gives $L_*q_M(\rho)$. Its fixed action is legal for every model. Known $\rho$ and compact support allow optimal centers; with an effective approximate codebook replace its distortion by its certified upper. The target is observed at load, so full-history service prediction is exact.
\end{proof}

The corrected predictable-start argument in~\eqref{eq:load-first-count} is essential. Subtracting an unweighted ``one last mark'' would give only an additive measure error and would not prove~\eqref{eq:load-occupation} for small sets.

\begin{corollary}[Two different raw geometric realizations]\label{cor:raw-geometries}
Theorem~\ref{thm:controlled-load} applies to each of the following families.

In a queue, take $W=1$ at preparation and, at service, transition to $W+1$ with probability $p_\theta(v,a)\le\bar p<1/2$ and to $W-1$ otherwise, ending at zero. Allow arbitrary finite noisy workload reports and all adaptive service actions. Then $L_*=(1-2\bar p)^{-1}$ is valid. An action-uniform exponential return envelope also holds.

In a finite quantum system, prepare any measurable state $\sigma_{\theta,v}$, including pure and rank-changing preparations. At each service use a legal normalized positive instrument whose alive part $\mathcal A_{\theta,v,a}$ satisfies
\[
 \mathcal A_{\theta,v,a}^*(I)\le(1-\zeta)I,
 \qquad 0<\zeta\le1.
\]
Then $L_*=1/\zeta$ is valid, uniformly over adaptive actions and retained states. Noncommuting instrument actions and preparations are allowed. In both families $\rho$ may be a Cantor product or a finite mixture of atoms and components of different dimension; the same actual concentration and cover bounds enter~\eqref{eq:load-risk}.
\end{corollary}
\begin{proof}
For the queue, stopped first-step drift of $W$ gives
$\E W_{n\wedge L}\le1-(1-2\bar p)\E(n\wedge L)$.
Nonnegativity and monotone convergence yield the displayed mean bound conditionally on every $v$ and prior history. The killed exponential test used in Theorem~\ref{thm:queue} is uniform in $v$ and gives the stronger tail. Since $p_\theta(v,a)$ can depend on the loaded mark and chosen action, occupation is genuinely length biased and policy dependent; it is not identified with $\rho$.

For the quantum family apply the dual inequality to the unnormalized alive state after each history-selected action. Induction bounds $\Pp(L>j\mid v,\text{past})\le(1-\zeta)^j$, and summing gives $\E L\le1/\zeta$. For an explicit noncommuting example on a qubit, let
\[
 K_{v,a}=U_a\diag(\sqrt{s_1(v)},\sqrt{s_2(v)}),\qquad
 B_v=\diag(\sqrt{1-s_1(v)},\sqrt{1-s_2(v)}),
\]
with $0<s_i(v)\le1-\zeta$, alive map $K_{v,a}\sigma K_{v,a}^*$ and end map $B_v\sigma B_v^*$. Their effects sum to $I$. Choose $U_0=I$ and $U_1$ the Hadamard unitary and take, for $v\in[0,1]$, $\sigma_v=(1-v)|0\rangle\langle0|+v|+\rangle\langle+|$. The preparations have rank two in the interior and rank one at the endpoints; the instrument actions do not share a state eigenbasis. State- and mark-dependent survival changes the acquired duration weights. No normalized nonlinear filter contraction is required for the held task factor, whose update after load is the identity. Its codebook is a global cover of $K$, and its error does not accumulate. These facts verify the delayed-load, initialization, covering and stability assumptions used for the causal upper. For the compact-family temporal results, impose continuous primitive coefficient functions on a compact parameter set; the finite report kernels and the single fresh load with common reference $\rho$ then give finite-response continuity by Lemma~\ref{lem:response}. Mere measurability is enough for Theorem~\ref{thm:controlled-load}, but is not claimed to give that extra continuity. They do not identify the full queue or quantum predictive quotient with the mark coordinate.
\end{proof}
The first-service lower covers all controls obeying the primitive tail bound, not merely the fixed action used by the upper. A simulator for either realization still requires the complete scored interface, including duration and the environment audit, and its internal state is multiplied into the retained budget as in Section~\ref{sec:composition}.

\section{Common-task composition and a matched resource law}\label{sec:composition}
\subsection{General defect transfer}
We use $\TV(\mu,\nu)=\sup_A|\mu(A)-\nu(A)|$. A bounded $[0,1]$ score changes in expectation by at most this distance; a stochastic row changes in $\ell^1$ by at most twice it. Both conventions must not be used interchangeably.

\begin{proposition}[Task-sensitive causal comparison]\label{prop:defect}
Consider two experiments and a paired class of parameter-independent programs, including any simulator's private state. Assume both complete scored processes have task-compatible finite-call bridges with errors $\omega_N$ and $\omega'_N$. Suppose the simulator supplies all virtual reports before their target decisions, preserves the common action and audit interface, and compares the complete first $N$ scored physical calls with total-variation error at most $e_N$. Then
\begin{equation}\label{eq:task-defect}
 |g-g'|\le\inf_{N\ge1}\{\omega_N+e_N+\omega'_N\}.
\end{equation}
A task readout discrepancy $\zeta_N$ adds $\zeta_N$ to the braces. A per-call complete-kernel discrepancy $\eta$ gives $e_N\le\min(1,N\eta)$ by telescoping. These assertions require the stated bridge at both endpoints, on the joint predictor--simulator state. A one-way causal simulator gives only the corresponding one-way optimal-risk comparison unless a reverse simulator is also supplied.
\end{proposition}
\begin{proof}
Insert the two finite-call expected risks between the gains. Their difference is at most $e_N$ for a common $[0,1]$ score, with any declared readout discrepancy added. Each gain-to-finite-call difference is bounded at its own endpoint. Infimizing gives~\eqref{eq:task-defect}. Coupling until first discrepancy, or telescoping signed kernels for a complete marked process, gives the per-call bound. A simulator using $s$ labels and a target predictor using $m$ labels is implemented by their joint $sm$-label state, with its actual clock and physical-call schedule. Evaluating recurrence or task compatibility only on the predictor coordinate would omit a required hypothesis.
\end{proof}
This is a clock-aligned comparison. Different physical clock maps require their own duration normalization; equality of a virtual call count does not make physical-time gains equal.

For quantum primitive comparisons the telescoping is performed on unnormalized marked completely positive maps in trace norm, not by coupling noncommuting states or dividing by a zero-probability report. A supplied Kraus-coordinate approximation can imply an instrument-norm certificate, but the resulting operational bound is independent of the Kraus representation. Report marginals without the duration, action and audit interface do not establish~\eqref{eq:task-defect}.

Equation~\eqref{eq:task-defect} allows amplification weaker than a uniform group-inverse estimate. For instance if both bridges are $O(N^{-a})$ and $e_N\le N\eta$, it gives $O(\eta^{a/(a+1)})$ by balancing. This is an upper modulus, not an automatically attained lower bound. The next construction supplies actual alternatives and a matching law rather than treating an allowed error tolerance as an irreducible loss.

\subsection{A single fully charged scored protocol}
Let $V$ have a fixed preparation law $\rho$ on a compact $K\subset\R^q$, independently in every round and independently of $\theta$. Let $\sigma\in\{-1,1\}$ be a fixed inaccessible calibration sign, not redrawn along the path. Let $0\le\delta\le1$ and $0\le e\le1/2$ be known magnitudes. Each round has four physical calls with visible types:
\begin{enumerate}
\item prepare fresh physical diagnostic data and a fresh mark $V$, retaining the observer label;
\item choose the diagnostic decision, score it, and then reveal the diagnostic report with intensity $\lambda_e(t)=(t^k-e)_+$;
\item report $V$ and update the retained label;
\item produce the fixed label readout $(c,b)\in\conv K\times[-1,1]$ before the audit and score $\norm{V-c}^2+(b-\sigma\delta)^2$.
\end{enumerate}
The diagnostic score is $t\one_{\{u\ne s\}}$ as before. The environment-owned audits in calls 2 and 4 are never feedback. Their alphabets and the early diagnostic-report deadline are part of the common measurable task interface. A simulator cannot replace the ground-truth audit by a fabricated one. All preparation, load and audit calls are charged; the physical acquisition budget is $N=4n$. Define $\calR$ to be expected total score divided by $n$ rounds. The usual physical-call average is $\calR/4$; dividing further by $D_K^2+5$ gives a $[0,1]$ normalization for each round's total score. Neither rescaling changes the resource exponents.

\begin{theorem}[Matched acquisition--memory--calibration--defect law]\label{thm:joint}
For every $n\ge1$, $M\ge2$, $0\le\delta\le1$ and $0\le e\le1/2$, the minimax value over all legal retained $M$-label programs in this protocol satisfies
\begin{equation}\label{eq:joint-sandwich}
 q_M(\rho)+\delta^2+D_n(\lambda_e)
 \le\calR(4n,M,\delta,e)
 \le q_{\lfloor M/2\rfloor}(\rho)+\delta^2+D_n(\lambda_e).
\end{equation}
The upper is realized by one stationary program using $2\lfloor M/2\rfloor$ labels. The lower diagnostic term holds even against full-history strategies. Uniformly over the displayed range,
\begin{equation}\label{eq:diagnostic-erasure-rate}
 D_n(\lambda_e)\asymp_k n^{-1/k}+e^{1/k}.
\end{equation}
Let $\varepsilon_*(e)$ denote source-to-target deficiency, uniformly in $\theta$, for one diagnostic interface with source intensity $\lambda_e$, target intensity $t^k$, the unchanged environment-owned audit, and the same early-report deadline. Then
\begin{equation}\label{eq:actual-deficiency}
 e/2\le\varepsilon_*(e)\le e.
\end{equation}
If the actual mark law has matching global covering and positive-mass small-ball dimension $d$, then
\begin{equation}\label{eq:joint-law}
 \calR(N,M,\delta,e)
 \asymp M^{-2/d}+N^{-1/k}+\delta^2+\varepsilon_*(e)^{1/k},\qquad N=4n.
\end{equation}
Constants depend on the displayed raw geometric data and $k$, not on $N,M,\delta,e$. No ambient dimension is substituted for that acquired dimension.
\end{theorem}
\begin{proof}
A program has at most $M$ geometric readouts. For every fresh $V$, irrespective of its earlier retained information, its geometric loss is at least the squared distance to that codebook; averaging gives $q_M(\rho)$. Average the two fixed calibration models $\sigma=\pm1$. Since their audit is never feedback and their visible experiments agree, the expected calibration loss is $\E b^2+\delta^2\ge\delta^2$. Independently average the two diagnostic models $\theta=\pm t$. Fresh geometric reports do not reveal their sign. The indistinguishable-before-revelation argument in Theorem~\ref{thm:diagnostic} therefore gives $D_n(\lambda_e)$. These lower terms hold simultaneously under the product of the two binary minimax priors, with $t$ chosen arbitrarily close to its supremum. They may therefore be added before taking the worst model.

For the upper, store a sign label and one of $J=\lfloor M/2\rfloor$ nearest geometric centers. Initialize the sign fairly. Retain it on $\bot$, update it on conclusive reports, and update only the center coordinate on a fresh load. Set $b=0$. The known physical protocol types select the appropriate row; no internal round count is provided. This uses $2J\le M$ labels and gives exactly the upper expression, or arbitrarily closely if a codebook infimum were not attained. Here compact finite-dimensional support does give attainment. Unused labels are assigned the same legal updates as a used label; no extra inaccessible information is stored there.

For~\eqref{eq:diagnostic-erasure-rate}, choosing $t=n^{-1/k}$ gives the $n^{-1/k}$ lower because $\lambda_e(t)\le1/n$; choosing $t=e^{1/k}$ gives the $e^{1/k}/2$ lower because no report is ever conclusive. The maximum of the two is at least half their sum. For the upper, if $t^k\le2e$ the value is at most $t/2\le(2e)^{1/k}/2$. Otherwise $\lambda_e(t)\ge t^k/2$ and the geometric sum is bounded by $\min(n,2/t^k)$, giving a constant times $n^{-1/k}$. This proves the uniform comparison, including $e=0$.

The identity simulator has complete report discrepancy $t^k-(t^k-e)_+\le e$, proving the upper in~\eqref{eq:actual-deficiency}. At $t=e^{1/k}$ the source produces only $\bot$ under both signs. Before the required target report the simulator has the same information in these two models. The two target report laws are $\mu_+=(1-e)\delta_\bot+e\delta_+$ and $\mu_-=(1-e)\delta_\bot+e\delta_-$; their total-variation distance is $e$. Any common simulated law has distance at least $e/2$ from one by the triangle inequality. Additional independent mark data and unchanged audits unavailable before the deadline do not alter this argument. Finally Lemma~\ref{lem:quantization}, the global cover, and $\lfloor M/2\rfloor\ge M/3$ for $M\ge2$ compare the two quantization terms and prove~\eqref{eq:joint-law}.
\end{proof}
For $e>0$ the blind threshold can make the individual diagnostic gain discontinuous. The sharp law above is obtained by a direct minimax argument; it does \emph{not} apply Proposition~\ref{prop:defect} while silently dropping an endpoint regularity condition. The quantity $\varepsilon_*$ is the deficiency of a real pair of raw interfaces and is bounded, not falsely asserted to equal the allowed symbol $e$.

\begin{corollary}[Singular acquired marks]\label{cor:singular-marks}
Theorem~\ref{thm:joint} applies to the standard middle-third Cantor preparation, with $d=\log2/\log3$, and to products of such measures and interval measures, with their summed acquired dimensions. More generally a finite mixture with positive weight on a highest-dimensional Ahlfors-regular component and global covering of that order has the same exponent, retaining that component's actual weight in the lower constant.
\end{corollary}
\begin{proof}
A depth-$j$ Cantor cylinder has mass $2^{-j}$ and diameter $3^{-j}$. Choosing $j$ so that $3^{-j}$ is comparable to $r$ shows that an interval of radius $r$ intersects at most a fixed number of depth-$j$ cylinders, and hence has mass at most $Cr^{\log2/\log3}$. The $2^j$ cylinders supply a global cover at that scale. Products preserve these estimates by enclosing Euclidean balls in product intervals and using product covers. For a finite mixture, keep the highest-dimensional component as the submeasure in~\eqref{eq:smallball}; cover the union, using that lower-dimensional components require no larger order. These are actual preparation and scored occupation measures of the four-call protocol, not conditional laws obtained by discarding low-probability outcomes.
\end{proof}

\subsection{Separate computational resources}
If geometric centers and readouts are represented to error $\xi$, the corresponding upper root distortion is at most $\sqrt{q_J(\rho)}+\xi$, so its squared term is at most $(\sqrt{q_J(\rho)}+\xi)^2$. In the general recursion the numerical term appears in~\eqref{eq:geometric-upper}. A particular finite implementation must supply terminating arithmetic or integration certificates for its row probabilities and decisions. A $p$-bit description of $J$ centers in $\R^q$ uses $O(Jqp)$ read-only bits, a nearest-center scan can use $O(qp+\log J)$ temporary bits, and it performs $O(Jq)$ fixed-precision distance operations per load. The retained sign--center coordinate uses $\lceil\log_2(2J)\rceil$ bits. These are explicit upper counts for the stated scan, not simultaneous optimal lower bounds on all resources. Continuous raw report reading and random-bit sampling remain part of the declared acquisition interface; their finite-precision implementation must be charged rather than assumed free.

\section{Sharp laws for general revelation profiles}\label{sec:profiles}
The convention at $\lambda(t)=0$ in each geometric sum is $\sum_{j=0}^{n-1}(1-\lambda(t))^j=n$. Every reference to all histories means legal visible histories and fresh private randomization, not future environment audit data.
The next result connects the general modulus to an all-strategy optimum; it is not used to prove Theorem~\ref{thm:modulus}. Let $a,\lambda:[0,1]\to[0,1]$ be continuous, with $a(0)=\lambda(0)=0$, $\lambda(t)>0$ for $t>0$, and $\lambda(1)=1$. A fixed hidden parameter is $\theta=st$, $s\in\{-1,1\}$. At each diagnostic call the observer selects a sign before the report, incurs the unobserved score $a(t)$ if wrong, and then receives the correct sign with probability $\lambda(t)$ and an inconclusive report otherwise. Reports are conditionally independent over calls. The score is never feedback. All model dependence resides in this positive raw kernel, not in the observer's program.

\begin{theorem}[A sharp task-visible acquisition profile]\label{thm:profile}
At every integer $n\ge1$, the minimax risk among all history-dependent strategies is
\begin{equation}\label{eq:profile-exact}
 D_n(a,\lambda)=\sup_{0\le t\le1}
       \frac{a(t)}{2n}\sum_{j=0}^{n-1}(1-\lambda(t))^j.
\end{equation}
One stationary two-label program attains this value at every horizon. Define
\begin{equation}\label{eq:profile-psi}
 \Psi_n(a,\lambda)=\sup_t a(t)
       \min\left\{1,\frac1{n\lambda(t)}\right\},
\end{equation}
with the minimum interpreted as one at $\lambda=0$. For the actual two-label marked family of this program,
\begin{equation}\label{eq:profile-modulus}
 K_n=\Psi_n(a,\lambda),\qquad
 \frac14\Psi_n\le D_n\le\frac12\Psi_n.
\end{equation}
The physical return time is one and $B_n=0$. Hence a continuum of raw experiments attains the general modulus against all histories, not only against a preselected controller.
\end{theorem}
\begin{proof}
At a fixed $t$, give the two signs equal prior probability. Until the first conclusive report, all visible histories and any private randomization have the same law for both signs, and the next sign choice has error probability $1/2$ under that prior. The probability of this event before call $j+1$ is $(1-\lambda(t))^j$. Later losses are nonnegative. This proves the lower bound for every history-dependent strategy and hence the minimax lower bound. A program initialized with a fair retained sign, unchanged on inconclusive reports and replaced on a conclusive report, has precisely that risk for each sign. It stores two labels and no clock, and so attains the lower bound simultaneously for all $n$ and $t$.

Relabel the two states as correct and wrong for mathematical evaluation under a fixed nonzero $t$. This relabeling is not information provided to the program. Its boundary arrays are
\[
 P_t=\begin{pmatrix}1&0\\\lambda(t)&1-\lambda(t)\end{pmatrix},
 \quad T_t=P_t,\quad r_t=(0,a(t))^T,\quad g_t=0.
\]
The same gain holds at zero since $a(0)=0$. A corrector with difference $v=u_{\rm wrong}-u_{\rm correct}$ has objective
$|a(t)-\lambda(t)v|+|v|/n$. Its minimum is
$a(t)\min(1,(n\lambda(t))^{-1})$: choose $v=0$ or $a(t)/\lambda(t)$. Conversely the vector
$z=(0,\min(1,(n\lambda(t))^{-1}))^T$
is feasible in~\eqref{eq:modulus-dual} and gives exactly that value. This proves the first identity in~\eqref{eq:profile-modulus}.

For $0<l\le1$, the geometric average is at most $\min(1,(nl)^{-1})$. It is at least half that quantity: $(1-l)^n\le(1+nl)^{-1}$ follows from $(1-l)^{-1}\ge1+l$ and the binomial inequality (with $l=1$ immediate), so
$[1-(1-l)^n]/(nl)\ge1/(1+nl)\ge\tfrac12\min(1,(nl)^{-1})$.
Multiply by $a(t)/2$, take the supremum and use the limiting convention at zero.
\end{proof}

Now introduce the same environment-owned fresh-mark and calibration task as in Theorem~\ref{thm:joint}, but replace its diagnostic amplitude $t$ and revelation probability $t^k$ by $a(t)$ and $\lambda(t)$. For clarity, a round has four charged calls: preparation, diagnostic decision and report, fresh-mark report, and terminal audit. The mark $V\sim\rho$ is independent and freshly prepared each round; its compact support $K$ is fixed. The audit scores
$\norm{V-c}^2+(b-\sigma\delta)^2$,
where $\sigma\in\{-1,1\}$ is fixed and never observed. A declared label has a fixed mark/calibration readout; the visible protocol phase is not a free remembered real number. Diagnostic and audit scores are not feedback. Risk in the next theorem is per round; division by four gives per physical call.

For $e\in[0,1/2]$ replace the diagnostic report probability by
$\lambda_e(t)=(\lambda(t)-e)_+$, with the same sign, action and unchanged audit wires. Define the visibility profile
\[
 A_*(u)=\sup\{a(t):\lambda(t)\le u\},\qquad 0\le u\le1.
\]
It tends to zero as $u\downarrow0$ by continuity and strict positivity of $\lambda$ away from zero. Suppose, when a simplified matching formula is desired, that
\begin{equation}\label{eq:profile-doubling}
 A_*(2u)\le C_a A_*(u)\quad(0<u\le1/2),\qquad
 q_{\lfloor M/2\rfloor}(\rho)\le C_q q_M(\rho)\quad(M\ge2).
\end{equation}
The exact sandwich below does not need these doubling conditions.

\begin{theorem}[A matched same-task multiresource law]\label{thm:profile-joint}
For every $n\ge1$, $M\ge2$, $0\le\delta\le1$, and $0\le e\le1/2$, the $M$-label minimax per-round risk satisfies
\begin{equation}\label{eq:profile-sandwich}
 q_M(\rho)+\delta^2+D_n(a,\lambda_e)
 \le\calR(4n,M,\delta,e)
 \le q_{\lfloor M/2\rfloor}(\rho)+\delta^2+D_n(a,\lambda_e).
\end{equation}
The source-to-target diagnostic deficiency $\varepsilon_*(e)$, with the virtual report due before any later action and without changing the environment audit, obeys
\begin{equation}\label{eq:profile-deficiency}
 e/2\le\varepsilon_*(e)\le e.
\end{equation}
Under~\eqref{eq:profile-doubling}, uniformly in the displayed budgets,
\begin{equation}\label{eq:profile-joint}
 \calR(4n,M,\delta,e)
 \asymp q_M(\rho)+\Psi_n(a,\lambda)+\delta^2+A_*(\varepsilon_*(e)).
\end{equation}
The upper program is parameter independent, uses $2\lfloor M/2\rfloor\le M$ labels, and is valid without an internally stored horizon. For a readout approximation of norm error at most $\xi$, its mark term is at most $(\sqrt{q_{\lfloor M/2\rfloor}(\rho)}+\xi)^2$. Program and workspace costs remain separate from this label count.
\end{theorem}
\begin{proof}
A fresh independent mark cannot be predicted from preceding rounds. Conditional on all earlier information, an $M$-state fixed-readout predictor has mark error at least $q_M(\rho)$; extra randomized readouts cannot improve the conditional squared error. Averaging the two signs of the invisible calibration shift gives a contribution at least $\delta^2$. At any fixed $t$, averaging the two diagnostic signs and applying the inconclusive-history argument of Theorem~\ref{thm:profile} gives $D_n$ before taking the supremum in $t$. These lower bounds are simultaneous under the product prior over the two independent signs, so they add. The loss is not a maximum of unrelated tasks.

For the upper bound retain the diagnostic sign together with a $J=\lfloor M/2\rfloor$-cell mark code. At the mark-report phase load the nearest-center index without changing the sign. At the audit output that center and zero for the calibration coordinate. Retain both labels through the next physical preparation. These rows give the upper risk in~\eqref{eq:profile-sandwich}; all phases are the visible four-call protocol. The numerical assertion follows from the $L^2$ triangle inequality.

The identity report map has error at most $e$ uniformly, giving the upper deficiency bound. By the intermediate value theorem choose $t$ with $\lambda(t)=e$. Both signs then have the same source report, whereas the two target report laws have total-variation distance $e$. Any common simulator output is at distance at least $e/2$ from one of them by the triangle inequality. It cannot consult the future audit. This proves~\eqref{eq:profile-deficiency}, not an unsupported exact equality with $e$.

Put $\Psi_{n,e}=\Psi_n(a,\lambda_e)$. Its definition gives
\[
 \Psi_{n,e}\ge\max\{\Psi_{n,0},A_*(e)\},\qquad
 \Psi_{n,e}\le\max\{2\Psi_{n,0},A_*(2e)\}.
\]
For the second inequality split $\lambda(t)>2e$, where $\lambda_e\ge\lambda/2$, from its complement. Thus $\Psi_{n,e}\asymp\Psi_{n,0}+A_*(e)$ under~\eqref{eq:profile-doubling}. The case $e=0$ follows directly. Equations~\eqref{eq:profile-deficiency} and~\eqref{eq:profile-doubling} compare $A_*(e)$ and $A_*(\varepsilon_*)$. The finite-horizon geometric-sum proof of Theorem~\ref{thm:profile} also applies to $\lambda_e$, even when it vanishes at positive $t$; its gain-zero and corrector identity are not asserted on those blind parameters. Use that finite-horizon estimate and quantization doubling to conclude.
\end{proof}

\begin{corollary}[Power and nonalgebraic rates]\label{cor:profile-examples}
For $a(t)=t^v$ and $\lambda(t)=t^k$, $v,k>0$,
\[
 \Psi_n=n^{-\min(v/k,1)},\qquad A_*(e)=e^{v/k}.
\]
For $\lambda(t)=t$ and $a(t)=1/\log(\mathrm e/t)$ on $t>0$, with $a(0)=0$,
\[
 \Psi_n=\frac1{\log(\mathrm e n)},\qquad
 A_*(e)=\frac1{\log(\mathrm e/e)}\quad(e>0).
\]
Both profiles satisfy the required doubling. Taking $\rho$ to be a Cantor product measure or a finite mixture with a positive highest-dimensional acquired component gives singular multiresource laws through its own $q_M(\rho)$.
\end{corollary}
\begin{proof}
For the powers split at $t=n^{-1/k}$. Above the split the expression is $t^{v-k}/n$, whose maximum is at the split if $v\le k$ and at one otherwise; below it $t^v$ is increasing. For the logarithmic example $a(t)$ increases and $a(t)/t=1/[t\log(\mathrm e/t)]$ decreases on $(0,1]$, so both pieces attain their maximum at $1/n$. The power doubling factor is $2^{v/k}$. In the logarithmic case it is at most $1+\log2$, since $e\le1/2$. The singular quantization assertions follow from the explicit prefix cover and retained submass proof of Theorem~\ref{thm:joint}; they do not turn a formal reachable set into an acquired distribution.
\end{proof}
The nonalgebraic example is deliberately outside the finite semialgebraic rate conclusion of Theorem~\ref{thm:algebraic}. Its exact two-label implementation and the converse in Theorem~\ref{thm:profile} remain valid. Accordingly a universal algebraic acquisition exponent cannot be inferred from compactness and task continuity alone.

\section{Scope of the general conclusion}\label{sec:scope}
The two-sided theorem concerns a finite retained boundary state and actual marked physical excursions, with all-row bounded mean duration. Its uniform limiting consequence additionally requires uniform integrability. The finite state belongs to the observer and the boundary interface, not necessarily to the physical experiment: the queue realization has unbounded physical state. No compactness is needed for the quantitative inequalities. Compactness and finite-response continuity enter the separate gain-continuity and robust-attainment theorem. There is no assertion of automatic positive-kernel compactness or unrestricted average-cost attainment.

The converse is for the maximal physical prefix discrepancy at the stated normalization, with an explicit return remainder. Terminal-time cancellation and within-cycle front loading rule out deleting these qualifications. The finite dual is an evaluation of an actual marked response, not a probability law and not a fast synthesis of every Borel program. The model-dependent analytical corrector and dual certificate cannot be supplied to a parameter-independent observer as free information.

The geometric statement uses actual scored occupation and its actual submass. Its finite-resolution lower requires only a concentration profile, but its causal upper still requires a common executable predictive initialization, global candidate-valid covers, report continuity and actual-law stability. An optimization over controllers also needs uniform lower witnesses for each competitor's own acquired law and an attaining fixed-size exploration, with the same finite-horizon baseline. These hypotheses are fully checked in the fresh-mark protocol; they are not asserted for unrestricted queue control or for arbitrary partially observed dynamics.

The profile theorem yields an exact all-history risk and a two-label implementation for arbitrary continuous diagnostic profiles. Its simplified joint resource law requires the displayed doubling conditions; without them the exact quantization sandwich is the conclusion. The erased blind-threshold family can have a nonzero or discontinuous limiting diagnostic gain, so its finite-horizon value is established directly and not passed through the gain-zero assertion or the compact uniform theorem. The calibration lower is for an actually invisible fixed sign, not an arbitrary error allowance.

The initial prepared-experiment problem is not exhausted by these results. Arbitrary nonregenerative histories, general nondominated common-program compactness, a necessary-and-sufficient spatial recursive geometry classification, and the optimal joint description, workspace, precision, physical-time and simulator-state region remain unproved. General data-dependent unbounded stopping and unbounded quantum operators require additional analysis. Total-variation error is used only for the specified scored interface and is not promoted to a large-deviation equivalence. The positive conclusion is a two-sided temporal obstruction with explicit physical accounting, coupled to finite-resolution acquired mass and to sharp implementable multiresource laws on verified raw experiments.

\begin{thebibliography}{99}
\bibitem{BasuPollackRoy}
S. Basu, R. Pollack and M.-F. Roy, \emph{Algorithms in Real Algebraic Geometry}, second edition, Algorithms and Computation in Mathematics 10, Springer, Berlin, 2006. doi:10.1007/3-540-33099-2.
\bibitem{BewleyKohlberg}
T. Bewley and E. Kohlberg, The asymptotic theory of stochastic games, \emph{Math. Oper. Res.} \textbf{1} (1976), 197--208. doi:10.1287/moor.1.3.197.
\bibitem{BewleyKohlbergStationary}
T. Bewley and E. Kohlberg, On stochastic games with stationary optimal strategies, \emph{Math. Oper. Res.} \textbf{3} (1978), 104--125. doi:10.1287/moor.3.2.104.
\bibitem{Blackwell}
D. Blackwell, Equivalent comparisons of experiments, \emph{Ann. Math. Statist.} \textbf{24} (1953), 265--272. doi:10.1214/aoms/1177729032.
\bibitem{ClarkeCheflesBarnettRiis}
R. B. M. Clarke, A. Chefles, S. M. Barnett and E. Riis, Experimental demonstration of optimal unambiguous state discrimination, \emph{Phys. Rev. A} \textbf{63} (2001), 040305. doi:10.1103/PhysRevA.63.040305.
\bibitem{DenardoFox}
E. V. Denardo and B. L. Fox, Multichain Markov renewal programs, \emph{SIAM J. Appl. Math.} \textbf{16} (1968), 468--487. doi:10.1137/0116038.
\bibitem{FedergrunSchweitzer}
A. Federgrun and P. J. Schweitzer, A fixed point approach to undiscounted Markov renewal programs, \emph{SIAM J. Algebraic Discrete Methods} \textbf{5} (1984), 539--550. doi:10.1137/0605052.
\bibitem{Fritz}
T. Fritz, A synthetic approach to Markov kernels, conditional independence and theorems on sufficient statistics, \emph{Adv. Math.} \textbf{370} (2020), 107239. doi:10.1016/j.aim.2020.107239.
\bibitem{GrafLuschgy}
S. Graf and H. Luschgy, \emph{Foundations of Quantization for Probability Distributions}, Lecture Notes in Mathematics 1730, Springer, Berlin, 2000. doi:10.1007/BFb0103945.
\bibitem{HellmanCover}
M. E. Hellman and T. M. Cover, Learning with finite memory, \emph{Ann. Math. Statist.} \textbf{41} (1970), 765--782. doi:10.1214/aoms/1177696958.
\bibitem{LamondPuterman}
B. F. Lamond and M. L. Puterman, Generalized inverses in discrete time Markov decision processes, \emph{SIAM J. Matrix Anal. Appl.} \textbf{10} (1989), 118--134. doi:10.1137/0610009.
\bibitem{Meyer}
C. D. Meyer, Jr., The role of the group generalized inverse in the theory of finite Markov chains, \emph{SIAM Rev.} \textbf{17} (1975), 443--464. doi:10.1137/1017044.
\bibitem{ShaliziCrutchfield}
C. R. Shalizi and J. P. Crutchfield, Computational mechanics: Pattern and prediction, structure and simplicity, \emph{J. Stat. Phys.} \textbf{104} (2001), 817--879. doi:10.1023/A:1010388907793.
\bibitem{Shaw}
S.-Y. Shaw, Convergence of ergodic nets and approximate solutions of linear functional equations, \emph{RIMS K\^oky\^uroku} \textbf{1415} (2005), 29--37.
\bibitem{GlynnInfanger}
P. W. Glynn and A. Infanger, Solution representations for Poisson's equation, martingale structure, and the Markov chain central limit theorem, \emph{Stochastic Systems} \textbf{14} (2024), 47--68; published online 28 December 2023. doi:10.1287/stsy.2022.0001.
\end{thebibliography}

\end{document}
